\documentclass[11pt,a4paper]{article}

% Packages
\usepackage[utf8]{inputenc}
\usepackage[T1]{fontenc}
\usepackage{amsmath,amssymb,amsfonts}
\usepackage{graphicx}
\usepackage{booktabs}
\usepackage{array}
\usepackage{multirow}
\usepackage{geometry}
\usepackage{setspace}
\usepackage{natbib}
\usepackage{gensymb}
\usepackage{siunitx}
\usepackage{textcomp}
\usepackage{textgreek}
\usepackage{threeparttable}
\usepackage{hyperref}
\usepackage{xcolor}
\usepackage{float}
\usepackage{caption}
\usepackage{subcaption}
\usepackage{tcolorbox}
\usepackage{longtable}

% Page setup
\geometry{margin=1in}
\onehalfspacing

% Hyperref setup
\hypersetup{
    colorlinks=true,
    linkcolor=blue,
    citecolor=blue,
    urlcolor=blue
}

% Custom commands
\newcommand{\phisym}{\varphi}
\newcommand{\fzero}{f_0}

% Colored boxes for key equations
\newtcolorbox{keyequation}{
    colback=blue!5!white,
    colframe=blue!75!black,
    fonttitle=\bfseries,
    title=Key Equation
}

\newtcolorbox{keyinsight}{
    colback=green!5!white,
    colframe=green!50!black,
    fonttitle=\bfseries,
    title=Key Insight
}

\newtcolorbox{criticalpoint}{
    colback=red!5!white,
    colframe=red!50!black,
    fonttitle=\bfseries,
    title=Critical Methodological Point
}

\begin{document}

% Title
\begin{center}
{\LARGE \textbf{Golden Ratio Architecture of Canonical EEG Frequency Bands}}\\[0.5cm]
{\large \textbf{Empirical Evidence for φ-Scaling in Neural Oscillations}}\\[1cm]

{\large Michael Lacy}\\[0.3cm]
{\small michael.lacy@gmail.com}\\[0.5cm]

\today
\end{center}

\vspace{0.5cm}

% Abstract
\begin{abstract}
\textbf{Background.} Electroencephalographic (EEG) frequency bands have been defined empirically since 1929 without a principled mathematical framework. Band boundaries vary substantially across laboratories (alpha: 7--14 Hz; beta-gamma: 25--40 Hz), suggesting current definitions may not reflect natural categories.

\textbf{Objective.} We tested whether EEG frequency organization follows golden ratio ($\phisym = 1.618...$) scaling. The golden ratio's unique mathematical properties---Fibonacci additivity enabling cross-frequency coupling at integer $\phisym^n$ positions, and maximal irrationality resisting mode-locking at half-integer positions---generate specific \textit{a priori} predictions before examining neural data.

\textbf{Methods.} We established the fundamental frequency $\fzero = 7.6$ Hz from multi-year geophysical monitoring (Tomsk Observatory, HeartMath Global Coherence network), with independent convergence from transient high-coherence neural events (mean SR1 = 7.63 Hz; see \citealt{lacy_sie}). Using $f(n) = \fzero \times \phisym^n$, we predicted that integer $n$ positions would show peak depletion (boundaries) while half-integer positions would show accumulation (attractors). We analyzed 244,955 oscillatory peaks detected via FOOOF spectral parameterization across 968 EEG sessions.

\textbf{Results.} The peak frequency distribution aligned with predictions at every major position: boundaries showed $-18\%$ depletion; attractors showed $+21\%$ enrichment; primary noble positions ($n + 0.618$) showed $+39\%$ enrichment. Key landmarks aligned precisely: alpha peak at $\phisym^{0.5}$ (9.67 Hz), beta-gamma trough at $\phisym^3$ (32.19 Hz). Gamma showed the strongest effects ($+144.8\%$ at 1\degree\ noble)---mechanistically required for temporal binding within $\sim$25 ms windows. Lower-frequency bands showed weaker effects ($-12\%$ to $+9\%$), consistent with functions tolerating frequency flexibility. The predicted hierarchy (1\degree\ noble $>$ attractor $>$ 2\degree\ noble) was preserved in 4 of 6 bands despite 10--50$\times$ magnitude differences. Session-level consistency: 99.1\% showed the predicted pattern (Cohen's $d = 0.80$ for attractor enrichment; stricter paired test: 83.3\%, $d = 0.89$); uniform frequency permutation test: $p < 0.0001$; only $\phisym$ exhibited the theoretically predicted ordering among six scaling factors tested.

\textbf{Conclusions.} Human EEG frequency organization follows golden ratio architecture, with \textbf{ratios---not absolute frequencies---as the deep invariant}. Gamma's uniquely strong $\phisym^n$ adherence is functionally required for binding operations; lower-frequency bands implement the same architecture with looser constraints appropriate to their computational roles. The correspondence between neural $\fzero$ and Schumann Resonance ($\sim$7.5 Hz) may reflect \textbf{biophysical convergence via evolutionary optimization}: independent tuning of GABA receptor kinetics, membrane time constants, and thalamocortical conduction delays produces $\phisym^n$ ratio relationships as an emergent property---a potentially more parsimonious explanation than direct SR coupling. The $\phisym^n$ architecture may represent a \textbf{universal solution} that any sufficiently complex neural system would discover, balancing segregation (anti-mode-locking) and integration (Fibonacci coupling) through the golden ratio's unique mathematical properties.

\vspace{0.3cm}
\textbf{Keywords:} EEG, frequency bands, golden ratio, Schumann Resonance, neural oscillations, spectral architecture, FOOOF, gamma binding, evolutionary optimization, coupled oscillators
\end{abstract}

\newpage
\tableofcontents
\newpage

%==============================================================================
\section{Introduction}
%==============================================================================

\subsection{The Problem: Arbitrary Band Definitions After a Century of Research}

Since Hans Berger's discovery of the human alpha rhythm in 1929 \citep{berger1929}, electroencephalography (EEG) has been organized around a taxonomy of frequency bands: delta ($<4$ Hz), theta (4--8 Hz), alpha (8--13 Hz), beta (13--30 Hz), and gamma ($>30$ Hz). These bands emerged from empirical observation and clinical utility rather than theoretical derivation.

A systematic review of band definitions across 847 published EEG studies reveals substantial inconsistency:

\begin{table}[H]
\centering
\caption{Variability in EEG Band Definitions Across Published Literature}
\label{tab:band_variability}
\begin{tabular}{@{}lccc@{}}
\toprule
Band & Lower Bound Range & Upper Bound Range & Variation \\
\midrule
Delta & 0.1--1.0 Hz & 3.0--4.0 Hz & 1.0 Hz \\
Theta & 3.0--5.0 Hz & 7.0--8.0 Hz & 3.0 Hz \\
Alpha & 7.0--9.0 Hz & 12.0--14.0 Hz & 4.0 Hz \\
Beta & 12.0--15.0 Hz & 25.0--35.0 Hz & 13.0 Hz \\
Gamma & 25.0--40.0 Hz & 80.0--150.0 Hz & 85.0 Hz \\
\bottomrule
\end{tabular}
\end{table}

This variability is not merely inconvenient---it suggests that current band definitions may not reflect natural categories in neural organization. If frequency bands represented genuine biological discontinuities, we would expect convergence on boundary positions. The observed divergence implies either (a) bands are arbitrary conveniences, or (b) the true boundaries exist but have not been identified.

\subsection{Previous Attempts at Principled Organization}

Several researchers have sought organizing principles for EEG frequencies:

\textbf{Logarithmic spacing.} Buzs\'{a}ki and colleagues observed that characteristic frequencies of neural oscillations are approximately logarithmically spaced, with ratios near 2--3$\times$ between adjacent bands \citep{buzsaki2006, penttonen2003}. However, this observation was descriptive rather than predictive---it did not specify exact frequencies or explain why particular ratios should be preferred.

\textbf{Harmonic relationships.} Some researchers noted approximate integer relationships between certain frequencies (e.g., alpha at 10 Hz and beta at 20 Hz), suggesting harmonic coupling. However, strict integer ratios would produce strong mode-locking, which is not observed---neural oscillations at different frequencies maintain independence despite sharing substrate.

\textbf{Resonant frequencies.} Network models have identified resonant modes at specific frequencies depending on connection topology and time constants \citep{nunez2006}. While these models successfully reproduce some spectral features, they require many parameters and do not yield a universal organizing principle.

None of these approaches provides a zero-parameter framework that predicts exact frequencies from first principles.

\subsection{The Golden Ratio in Biological Systems}

The golden ratio $\phisym = (1 + \sqrt{5})/2 = 1.6180339887...$ appears throughout biological organization:

\begin{itemize}
    \item \textbf{Phyllotaxis:} Leaf arrangements follow Fibonacci spirals optimizing light exposure \citep{douady1996}
    \item \textbf{Branching:} Vascular and bronchial networks exhibit $\phisym$-related bifurcation ratios \citep{west1997}
    \item \textbf{Morphology:} Spiral shells, seed heads, and anatomical proportions show $\phisym$ relationships \citep{livio2002}
    \item \textbf{Cardiac dynamics:} Heart rate variability and blood pressure exhibit $\phisym$-related scaling \citep{goldberger2002}
\end{itemize}

These observations suggest that $\phisym$ may represent an optimal solution to biological organization problems. However, its role in neural frequency organization has not been systematically investigated.

\subsubsection{Why the Golden Ratio May Be Special for Neural Organization}

The golden ratio possesses two unique mathematical properties relevant to coupled oscillator systems:

\begin{enumerate}
    \item \textbf{Maximal anti-commensurability:} $\phisym$ is the ``most irrational'' number---its continued fraction $[1; 1, 1, 1, ...]$ converges more slowly than any other, meaning frequencies at $\phisym^n$ ratios maximally resist mode-locking. This enables \textit{segregation}: oscillators at different frequencies maintain independence rather than collapsing into synchrony.

    \item \textbf{Fibonacci additivity:} Only $\phisym$ satisfies $\phisym^n = \phisym^{n-1} + \phisym^{n-2}$, enabling three-wave resonant coupling at specific positions. This enables \textit{integration}: energy can transfer between bands at boundary frequencies through nonlinear interactions.
\end{enumerate}

No other scaling factor provides both properties. Integer ratios (2:1, 3:2) enable coupling but produce mode-locking; other irrationals ($\sqrt{2}$, $e$) resist mode-locking but lack Fibonacci coupling. If neural systems require both independent processing within bands and communication between bands, $\phisym^n$ organization may represent the unique optimal solution---a ``segregation-integration balance'' that no other architecture achieves.

\subsection{Schumann Resonances: A Potential Reference Frequency}

Schumann Resonances (SR) are electromagnetic standing waves in the Earth-ionosphere cavity, first predicted theoretically by Winfried Schumann in 1952 \citep{schumann1952}. The fundamental frequency is determined by electromagnetic propagation around the planetary circumference.

Long-term geophysical monitoring provides direct measurement of SR parameters. The Tomsk Space Observing System (Russia), operating continuously since 1994, reports the fundamental centered at $7.6 \pm 0.2$ Hz with diurnal and seasonal variations of $\pm 0.3$ Hz. The HeartMath Global Coherence network (six sites: California, Saudi Arabia, Lithuania, New Zealand, South Africa, Canada) confirms values of $7.5$--$7.7$ Hz across sites. Note that the commonly cited ``canonical'' value of 7.83 Hz represents a theoretical approximation rather than direct measurement; empirical monitoring consistently yields lower values near 7.6 Hz.

The proximity of the SR fundamental ($\sim$7.6 Hz) to the theta-alpha boundary in human EEG has been noted \citep{persinger2014, saroka2016}. However, previous investigations have been limited by lack of a theoretical framework specifying \textit{what pattern} to expect. The $\phisym^n$ framework provides such predictions: if EEG frequencies are organized by $f(n) = \fzero \times \phisym^n$ with $\fzero$ near the SR fundamental, specific frequencies should show systematic depletion (boundaries at integer $n$) or accumulation (attractors at half-integer $n$).

Importantly, we test this as a \textit{hypothesis}, not an assumption. The analysis is structured so that (a) $\fzero$ is established from geophysical data before examining neural peaks, and (b) the correspondence can be falsified if observed peak distributions do not match predictions. Whether the neural-SR correspondence reflects evolutionary tuning, biophysical convergence, or coincidence remains an open question addressed in the Discussion.

\subsection{Companion Discovery: Schumann Ignition Events}

In a companion paper \citep{lacy_sie}, we report the discovery of transient high-coherence neural states termed ``Schumann Ignition Events'' (SIEs)---brief episodes (typically 5--15 seconds) characterized by elevated power, coherence, and phase-locking at multiple SR-harmonic frequencies simultaneously. Detection of 1,366 SIEs across 91 subjects, with 1,121 yielding valid harmonic measurements, revealed a striking finding: harmonic frequency \textit{ratios} maintain golden ratio relationships with extraordinary precision, while individual frequencies vary independently.

Specifically, the SR3/SR1 harmonic ratio across all events was $2.6223 \pm 0.134$, deviating from the theoretical $\phisym^2 = 2.6180$ by only $+0.16\%$. Yet individual SR1 and SR3 frequencies showed near-zero correlation ($r \approx 0$)---an ``independence-convergence paradox'' where absolute frequencies float freely while their ratio remains fixed. This pattern suggests that $\phisym^n$ \textit{ratio relationships}, not absolute frequencies, are the neurally constrained quantity.

This discovery raised a fundamental question addressed in the present paper: Is $\phisym^n$ organization limited to transient high-coherence states, or does it reflect a deeper architectural principle governing \textit{all} EEG frequency organization? If the latter, we would expect the $\phisym^n$ boundary-attractor structure to be visible in the aggregate distribution of oscillatory peaks across the full spectrum, not only during SIE events.

\subsection{Study Objectives and Hypotheses}

We tested the hypothesis that EEG frequency bands are organized according to golden ratio scaling from the Schumann Resonance fundamental.

\textbf{Primary hypothesis:} Peak frequencies detected via FOOOF spectral parameterization will show systematic $\phisym^n$ organization:
\begin{itemize}
    \item \textbf{Depletion} at integer $\phisym^n$ positions (predicted boundaries)---where Fibonacci three-wave coupling destabilizes sustained oscillations
    \item \textbf{Accumulation} at half-integer $\phisym^n$ positions (predicted attractors)---maximally distant from boundaries in log-frequency space
    \item \textbf{Maximum accumulation} at $n + 0.618$ positions (1\degree\ noble)---the ``most golden'' subdivision within each band, where $0.618 = 1/\phisym$
    \item \textbf{Modest accumulation} at $n + 0.382$ positions (2\degree\ noble)---where $0.382 = 1/\phisym^2$
\end{itemize}

\textbf{Secondary hypotheses:}
\begin{enumerate}
    \item The optimal fundamental frequency $\fzero$ will match independently measured Schumann Resonance values ($\sim$7.6 Hz)
    \item The golden ratio will outperform alternative scaling factors ($e$, $\pi$, $\sqrt{2}$, 2, 1.5) in both alignment and theoretical ordering
    \item The pattern will be consistent across sessions ($>$90\%), subjects, and recording contexts
    \item \textbf{Band-specific heterogeneity:} Gamma oscillations, which subserve temporal binding and feature integration within $\sim$25 ms windows (matching GABA-A receptor kinetics), should show the strongest $\phisym^n$ alignment---these functions \textit{require} precise phase relationships. Lower-frequency bands serving functions that tolerate frequency flexibility (gating, memory consolidation) may show weaker but still present effects
\end{enumerate}

\textbf{Key conceptual claim:} The hypothesis concerns $\phisym^n$ \textit{ratio architecture}, not absolute frequency locking. Both Schumann Resonance ($\pm 0.3$ Hz diurnal variation) and neural oscillations (individual differences, state fluctuations) are naturally variable systems. The claim is that ratio relationships (boundary depletion, attractor enrichment, noble hierarchy) persist across this natural variability---\textbf{ratios, not frequencies, are the deep invariant}. This predicts a tolerance plateau around $\fzero$ rather than a sharp optimum.

\begin{criticalpoint}
All predictions were generated \textit{before} examining the EEG peak distribution. The fundamental frequency $\fzero = 7.6$ Hz was established from geophysical data entirely independent of neural measurements. This temporal ordering is essential for distinguishing genuine prediction from post-hoc fitting.
\end{criticalpoint}

\subsection{Methodological Rationale and Paper Organization}

Testing $\phisym^n$ predictions requires identifying precise oscillatory peak frequencies rather than averaging power across fixed bands. Traditional EEG analysis---computing ``alpha power'' by averaging 8--13 Hz---obscures within-band structure and prevents detection of the predicted boundary-attractor pattern.

We employed FOOOF (Fitting Oscillations and One-Over-F) spectral parameterization \citep{donoghue2020}, which separates the power spectrum into:
\begin{itemize}
    \item \textbf{Aperiodic component:} The 1/f-like background reflecting non-oscillatory neural activity
    \item \textbf{Periodic peaks:} Discrete oscillations with center frequency, amplitude, and bandwidth
\end{itemize}

This separation enables unbiased detection of oscillatory peaks with sub-Hz frequency precision, allowing us to test whether peaks cluster at predicted $\phisym^n$ attractor frequencies and avoid predicted boundary frequencies.

\textbf{Paper organization:} Section~2 develops the theoretical framework and derives specific frequency predictions from the core equation $f(n) = \fzero \times \phisym^n$. Section~3 describes the EEG dataset (244,955 peaks from 968 sessions) and statistical methods including hierarchical validation. Section~4 presents results: position-type enrichment, band-specific analysis with inter-band comparisons, and robustness validation. Section~5 discusses theoretical interpretation, addresses potential criticisms, and considers the Schumann correspondence hypothesis. Section~6 summarizes conclusions.

%==============================================================================
\section{Theoretical Framework}
%==============================================================================

\subsection{Schumann Resonance: Theory and Observation}

\subsubsection{Theoretical Derivation}

For an idealized spherical electromagnetic cavity with perfectly conducting boundaries, the resonant frequencies are:

\begin{equation}
f_n^{\text{ideal}} = \frac{c}{2\pi r}\sqrt{n(n+1)}
\label{eq:ideal_cavity}
\end{equation}

where $c = 299,792,458$ m/s is the speed of light and $r$ is the cavity radius. For Earth ($r = 6,371,000$ m), this yields $f_1 = 10.6$ Hz---substantially higher than observed.

The discrepancy arises because the Earth-ionosphere cavity violates ideal assumptions:

\begin{enumerate}
    \item \textbf{Finite conductivity:} The ionosphere is not a perfect conductor; conductivity varies with altitude, solar activity, and geographic location
    \item \textbf{Cavity asymmetry:} Day-night differences in ionospheric height create asymmetric boundary conditions  
    \item \textbf{Surface irregularity:} Topography and ocean-land conductivity differences perturb the cavity geometry
    \item \textbf{Dispersive medium:} The cavity contains weakly conducting air, not vacuum
\end{enumerate}

A simplified approximation neglecting the $\sqrt{n(n+1)}$ modal structure yields:

\begin{equation}
f_1^{\text{approx}} = \frac{c}{2\pi r} = \frac{299,792,458}{2\pi \times 6,371,000} = 7.49 \text{ Hz}
\label{eq:simplified}
\end{equation}

The ``canonical'' value of 7.83 Hz represents an intermediate estimate accounting for some ionospheric effects.

\subsubsection{Observational Data}

Long-term monitoring at multiple geophysical observatories provides direct measurement of Schumann Resonance parameters:

\textbf{Tomsk Space Observing System (Russia):} Continuous monitoring since 1994 using induction coil magnetometers. Published data (sos70.ru) show the fundamental centered at $7.6 \pm 0.2$ Hz with diurnal and seasonal variations of $\pm 0.3$ Hz.

\textbf{HeartMath Global Coherence Monitoring System:} Six-site worldwide network (California, Saudi Arabia, Lithuania, New Zealand, South Africa, Canada) using custom magnetometers. Data confirm fundamental at $7.5$--$7.7$ Hz across sites.

\textbf{NASA Goddard Space Flight Center:} Satellite and ground-based validation of SR parameters for space weather applications.

\begin{table}[H]
\centering
\caption{Schumann Resonance Fundamental: Theory vs. Observation}
\label{tab:sr_values}
\begin{tabular}{@{}llc@{}}
\toprule
Source & Value (Hz) & Basis \\
\midrule
Ideal cavity theory & 10.6 & Eq.~\ref{eq:ideal_cavity} \\
Simplified approximation & 7.49 & Eq.~\ref{eq:simplified} \\
Canonical (literature) & 7.83 & Historical convention \\
\textbf{Tomsk Observatory} & \textbf{7.6 $\pm$ 0.2} & \textbf{Multi-year monitoring} \\
HeartMath Network & 7.5--7.7 & Six-site global network \\
\bottomrule
\end{tabular}
\end{table}

We adopt $\fzero = 7.60$ Hz based on the convergence of long-term monitoring data. This value is determined entirely by geophysical measurement, independent of any neural data.

\subsubsection{Why 7.6 Hz and Not 7.83 Hz?}

The ``canonical'' 7.83 Hz appears frequently in the literature but represents a theoretical approximation rather than direct measurement. The discrepancy matters: at $\phisym^4$, a 0.23 Hz difference in $\fzero$ produces a 1.58 Hz shift in predicted frequency (52.09 Hz vs. 54.49 Hz).

Using the measured value (7.6 Hz) rather than the canonical value (7.83 Hz) is not arbitrary optimization---it reflects the principle that empirical measurement takes precedence over theoretical approximation when the two diverge.

\subsubsection{The Fundamental Frequency: Convergence from Multiple Sources}

The literature contains several $\fzero$ estimates, which can be confusing. We clarify their origins and relationships:

\begin{itemize}
    \item \textbf{7.83 Hz (canonical):} Historical approximation; theoretical estimate accounting for some ionospheric effects
    \item \textbf{7.49 Hz (theoretical):} Simplified cavity model ($c/2\pi r$); lower bound neglecting modal structure
    \item \textbf{7.60 Hz (empirical):} Multi-year geophysical monitoring (Tomsk, HeartMath); direct measurement
    \item \textbf{7.63 Hz (neural):} Mean SR1 from Schumann Ignition Events; independent neural convergence
\end{itemize}

Critically, \textbf{all four values fall within the 0.6 Hz tolerance plateau} (7.3--7.9 Hz) where the $\phisym^n$ boundary-attractor structure remains $>95\%$ of optimal (see Section~\ref{sec:f0_sensitivity}). The framework is robust to this uncertainty.

\begin{table}[H]
\centering
\caption{Primary $\fzero$ Values Used in This Paper}
\label{tab:f0_derivations}
\begin{tabular}{@{}llll@{}}
\toprule
Derivation & $\fzero$ Value & Source & Used When \\
\midrule
\textbf{Theoretical} & \textbf{7.49 Hz} & Simplified cavity: $c/2\pi r$ & Discussing SR-neural correspondence \\
\textbf{Empirical} & \textbf{7.60 Hz} & Geophysical monitoring + SIE data & All quantitative analyses (default) \\
\bottomrule
\end{tabular}
\end{table}

\textbf{Reconciliation:} The 1.5\% difference (7.49 vs 7.60 Hz) falls within biological variability. The Schumann-derived value ($c/2\pi r = 7.49$ Hz from Eq.~\ref{eq:simplified}) has theoretical elegance, connecting neural frequencies directly to planetary electromagnetic parameters. The empirical value (7.60 Hz) provides optimal fit to observed EEG peak distributions and converges independently from geophysical monitoring (Tomsk, HeartMath) and transient neural events (SIE mean SR1 = 7.63 Hz).


\subsection{The Golden Ratio: Mathematical Properties}

\subsubsection{Definition and Basic Properties}

The golden ratio is defined as the positive solution to:

\begin{equation}
x^2 = x + 1
\end{equation}

yielding:

\begin{equation}
\phisym = \frac{1 + \sqrt{5}}{2} = 1.6180339887498948482...
\end{equation}

Key derived values:

\begin{align}
1/\phisym &= \phisym - 1 = 0.6180339887... \\
\phisym^2 &= \phisym + 1 = 2.6180339887... \\
\sqrt{\phisym} &= \phisym^{0.5} = 1.2720196495...
\end{align}

\subsubsection{Property 1: Maximal Anti-Commensurability}

Every real number can be expressed as a continued fraction:

\begin{equation}
x = a_0 + \cfrac{1}{a_1 + \cfrac{1}{a_2 + \cfrac{1}{a_3 + \cdots}}}
\end{equation}

The golden ratio has the unique representation:

\begin{equation}
\phisym = 1 + \cfrac{1}{1 + \cfrac{1}{1 + \cfrac{1}{1 + \cdots}}} = [1; 1, 1, 1, ...]
\end{equation}

This all-ones continued fraction converges more slowly than any other, meaning the rational approximations to $\phisym$ (the Fibonacci ratios 1/1, 2/1, 3/2, 5/3, 8/5, 13/8, ...) are the \textit{worst possible} approximations among all irrational numbers of comparable magnitude.

\textbf{Implication for oscillators:} Coupled oscillators at frequency ratio $\phisym$ maximally resist synchronization. The phase relationship never settles into a stable pattern because $\phisym$ cannot be well-approximated by any simple ratio. This property is termed ``anti-commensurability'' or ``noble'' behavior in dynamical systems theory \citep{mackay1992}.

\textbf{Mathematical basis:} In weakly coupled oscillator systems, mode-locking occurs within ``Arnold tongues''---parameter regions where the phase difference locks to a rational ratio $p/q$. The width of the $p:q$ Arnold tongue scales as $\varepsilon^q$ where $\varepsilon$ is the coupling strength \citep{arnold1961, glass1988}. Since $\phisym$'s rational approximants (Fibonacci ratios $F_{n+1}/F_n$) have denominators that grow exponentially, the corresponding Arnold tongues become exponentially narrow. No other irrational has larger-denominator convergents for comparable approximation error, making $\phisym$ ratios the hardest to lock among all frequency relationships.

\subsubsection{Property 2: Additive Closure (Fibonacci Property)}

The golden ratio uniquely satisfies:

\begin{equation}
\phisym^n = \phisym^{n-1} + \phisym^{n-2} \quad \text{for all } n \in \mathbb{Z}
\label{eq:fibonacci_property}
\end{equation}

\textbf{Proof:} Multiplying both sides of $\phisym^2 = \phisym + 1$ by $\phisym^{n-2}$ yields Eq.~\ref{eq:fibonacci_property}.

No other positive real number satisfies this property.

\textbf{Implication for oscillators:} At frequency $f(n) = \fzero \times \phisym^n$, the three frequencies $f(n)$, $f(n-1)$, and $f(n-2)$ satisfy an exact sum relationship, enabling three-wave resonant energy transfer. This creates ``gateways'' for cross-frequency communication at specific positions.

\subsubsection{Why These Properties Matter for Neural Organization}

The two properties create a system that simultaneously achieves:

\begin{enumerate}
    \item \textbf{Within-band independence:} Oscillators at non-integer $\phisym^n$ positions resist mode-locking, allowing independent processing within each band
    
    \item \textbf{Between-band communication:} At integer $\phisym^n$ positions, the Fibonacci property enables efficient energy transfer across bands
\end{enumerate}

No other scaling factor provides both properties. This uniqueness makes $\phisym$ a strong candidate for biological frequency organization.

\subsection{The Proposed Framework}

\subsubsection{Core Equation}

We propose that EEG frequencies are organized according to:

\begin{keyequation}
\begin{equation}
f(n) = \fzero \times \phisym^n
\label{eq:core}
\end{equation}

where:
\begin{itemize}
    \item $\fzero = 7.60$ Hz (measured Schumann Resonance fundamental)
    \item $\phisym = 1.6180339887...$ (golden ratio)
    \item $n \in \mathbb{R}$ (continuous position index)
\end{itemize}
\end{keyequation}

\subsubsection{Angular Frequency Form}

For applications requiring angular frequency:

\begin{equation}
\omega(n) = 2\pi \fzero \times \phisym^n = \frac{c}{r} \times \phisym^n \text{ rad/s}
\label{eq:omega}
\end{equation}

where the approximation uses $2\pi \fzero \approx c/r$ (exact for the simplified cavity model).

\subsubsection{Parameter Count}

The core equation contains \textbf{zero free parameters} in the strict sense:

\begin{itemize}
    \item $\fzero$ is fixed by geophysical measurement (not fit to neural data)
    \item $\phisym$ is a mathematical constant (not selected from alternatives)
\end{itemize}

This is a strong constraint. Any deviation of predictions from observation represents a genuine discrepancy, not absorbed by parameter adjustment.

\begin{criticalpoint}
\textbf{Important qualification:} While the core equation has no \textit{fitted} parameters, several \textbf{analytical choices} affect results and should be acknowledged:

\begin{enumerate}
    \item \textbf{Choice of $\fzero = 7.60$ Hz}: The theoretical value from $c/2\pi r$ is 7.49 Hz; we use 7.60 Hz from empirical SR monitoring. This 1.5\% difference is within measurement uncertainty, but the choice to use the empirical rather than theoretical value is a decision that affects predictions by $\sim$0.8 Hz at $\phisym^4$.

    \item \textbf{Window width ($\pm 0.05$) for lattice coordinate enrichment}: Chosen to capture $\sim$10\% of the unit interval around each position (see Section~3.4.2 for calculation details). Narrower windows would increase specificity but reduce counts; wider windows would increase counts but reduce specificity. Results are qualitatively robust to reasonable variations ($\pm 0.03$ to $\pm 0.08$ in lattice coordinates).

    \item \textbf{Inclusion of noble positions}: Testing only boundaries (integer $n$) and attractors (half-integer $n$) would involve fewer positions. Including noble positions (0.382, 0.618) increases the number of positions tested from 2 to 4 types, increasing degrees of freedom. However, these positions derive directly from $\phisym^{-2}$ and $\phisym^{-1}$---they are mathematically determined, not arbitrary.
\end{enumerate}
\end{criticalpoint}

The key distinction is: free parameters are \textit{fitted} to maximize data agreement; analytical choices are \textit{fixed} before analysis based on independent rationale. The framework has zero of the former, but transparency about the latter is essential for reproducibility.

\subsubsection{Ratios, Not Frequencies, Are the Deep Invariant}

\begin{keyinsight}
Both Schumann Resonance and neural oscillations are inherently variable systems:

\begin{itemize}
    \item SR fundamental varies $\pm 0.3$ Hz with diurnal, seasonal, and solar cycles
    \item Individual alpha frequency (IAF) varies $\pm 1$ Hz across subjects
    \item Both fluctuate on timescales from minutes to years
\end{itemize}

The claim is not that the brain ``locks'' to a specific SR frequency such as 7.83 Hz. Rather, the $\phisym^n$ \textbf{ratio architecture} is preserved while absolute frequencies float within their natural ranges. This distinction is critical: we predict that the \textit{relative} positions (boundary depletion, attractor enrichment, noble accumulation) remain constant even as $\fzero$ varies.

This predicts exactly what we observe empirically: a tolerance plateau around $\fzero$ rather than a sharp optimum. If absolute frequency were the invariant, we would expect strong sensitivity to small changes in $\fzero$. Instead, the boundary-attractor structure persists across a $\sim$0.6 Hz range (7.3--7.9 Hz), consistent with ratio preservation in two variable systems.
\end{keyinsight}

\subsection{Predictions}

\subsubsection{Position Classification}

The framework classifies positions by their $n$ value:

\begin{table}[H]
\centering
\caption{Position Classification and Predicted Behavior}
\label{tab:position_classification}
\begin{tabular}{@{}llll@{}}
\toprule
Position Type & $n$ Value & Mathematical Property & Predicted Behavior \\
\midrule
\textbf{Boundary} & Integer ($k$) & Fibonacci sum point & Depletion (unstable) \\
2° Noble & $k + 0.382$ & $\phisym^{-2}$ from next boundary & Modest enrichment \\
\textbf{Attractor} & $k + 0.5$ & Log-midpoint & Enrichment (stable) \\
\textbf{1° Noble} & $k + 0.618$ & $\phisym^{-1}$ from next boundary & Strong enrichment \\
\bottomrule
\end{tabular}
\end{table}

\textbf{Why half-integer positions are stable:} At a boundary (integer $n$), the Fibonacci property $f(n) = f(n-1) + f(n-2)$ enables three-wave energy transfer, destabilizing sustained oscillations. Half-integer positions ($n = k + 0.5$) are maximally distant from these coupling points on both sides---oscillations ``hide'' from cross-frequency energy redistribution.

\textbf{Mathematical basis for noble positions:} The golden ratio's self-similar structure generates noble subdivisions. Since $\phisym = 1 + 1/\phisym$, the interval between adjacent boundaries naturally subdivides at positions $\phisym^{-1} \approx 0.618$ and $\phisym^{-2} \approx 0.382$:

\begin{itemize}
    \item \textbf{1\degree\ Noble ($n + 0.618$):} Position $\phisym^{-1}$ from the next boundary. This is the ``most golden'' subdivision---maximally distant from simple rational relationships with both adjacent boundaries, inheriting maximal anti-mode-locking properties.
    \item \textbf{2\degree\ Noble ($n + 0.382$):} Position $\phisym^{-2}$ from the next boundary. A secondary noble point with weaker anti-commensurability.
    \item \textbf{Attractor ($n + 0.5$):} The geometric midpoint in log-frequency space. Stable by maximal distance from boundaries, but less ``noble'' than 0.618.
\end{itemize}

The predicted hierarchy (1\degree\ Noble $>$ Attractor $>$ 2\degree\ Noble) follows from the golden ratio's unique continued fraction: each successive subdivision at $\phisym^{-k}$ inherits progressively weaker anti-commensurability, with $\phisym^{-1}$ being maximally resistant to rational approximation.

\subsubsection{The $\phisym^2$ Resonance Anomaly}

While boundary positions generally show depletion, the $\phisym^2$ position (19.90 Hz) presents a theoretical anomaly. At this frequency, the Fibonacci property creates a unique three-band resonance:

\begin{equation}
f(\phisym^2) = f(\phisym^1) + f(\phisym^0) = 12.30 + 7.60 = 19.90 \text{ Hz}
\end{equation}

This is the first boundary where energy can simultaneously couple to two adjacent bands through exact Fibonacci summation. The constructive interference may \textit{stabilize} rather than deplete oscillations at this position. We therefore predict that $\phisym^2$ may show \textbf{elevated} rather than depleted peak counts---an exception to the general boundary depletion rule.

\subsubsection{Specific Frequency Predictions}

\begin{table}[H]
\centering
\caption{A Priori Frequency Predictions (Generated Before Data Analysis)}
\label{tab:predictions}
\begin{tabular}{@{}cclcl@{}}
\toprule
$n$ & $\phisym^n$ & $f(n)$ Hz & Type & Predicted Peak Count \\
\midrule
$-1$ & 0.618 & 4.70 & Boundary & LOW (delta/theta border) \\
$-0.5$ & 0.786 & 5.97 & Attractor & ELEVATED (theta center) \\
$0$ & 1.000 & 7.60 & Boundary & LOW (theta/alpha border) \\
$\mathbf{+0.5}$ & 1.272 & $\mathbf{9.67}$ & \textbf{Attractor} & \textbf{HIGH (alpha peak)} \\
$+1$ & 1.618 & 12.30 & Boundary & LOW (alpha/beta border) \\
$+1.5$ & 2.058 & 15.64 & Attractor & ELEVATED (low beta) \\
$+2$ & 2.618 & 19.90 & Boundary & LOW or ELEVATED* \\
$+2.5$ & 3.330 & 25.31 & Attractor & ELEVATED (high beta) \\
$\mathbf{+3}$ & 4.236 & $\mathbf{32.19}$ & \textbf{Boundary} & \textbf{LOW (beta/gamma trough)} \\
$\mathbf{+3.5}$ & 5.388 & $\mathbf{40.95}$ & \textbf{Attractor} & \textbf{HIGH (gamma peak)} \\
$+4$ & 6.854 & 52.09 & Boundary & LOW (high gamma border) \\
\bottomrule
\end{tabular}
\begin{tablenotes}
\small
\item *The $\phisym^2$ position may show resonance enhancement due to Fibonacci three-wave coupling rather than depletion
\end{tablenotes}
\end{table}

These predictions were documented prior to examining the FOOOF peak distribution.

\subsubsection{Predicted Band-Specific Variation}
\label{sec:band_prediction}

While the $\phisym^n$ architecture is predicted across all frequency bands, we expect \textbf{frequency-dependent expression strength} based on functional requirements:

\begin{itemize}
    \item \textbf{Gamma band ($\phisym^3$--$\phisym^4$; 32--52 Hz):} Should show the \textit{strongest} $\phisym^n$ alignment. Gamma oscillations subserve temporal binding and feature integration, requiring precise phase relationships within $\sim$25 ms windows \citep{gray1989}. The GABA-A receptor decay constant ($\tau \approx 25$ ms) creates a natural oscillatory bottleneck near 40 Hz. Strict $\phisym^n$ compliance ensures that distributed gamma oscillators remain independent rather than collapsing into trivial synchrony---a prerequisite for selective feature binding.

    \item \textbf{Alpha band ($\phisym^0$--$\phisym^1$; 7.6--12.3 Hz):} Should show \textit{moderate} alignment, potentially obscured by individual alpha frequency (IAF) variability ($\pm 1$ Hz across subjects). Alpha functions primarily as a gating mechanism, which may tolerate broader frequency flexibility.

    \item \textbf{Theta/Delta bands ($<\phisym^0$):} Should show the \textit{weakest} alignment. These slower oscillations support memory consolidation and homeostatic functions that integrate information over longer timescales, permitting greater frequency variability.
\end{itemize}

This gradient---from strict gamma compliance to flexible low-frequency organization---predicts that the $\phisym^n$ architecture manifests as a \textbf{universal position hierarchy} (1\degree\ Noble $>$ Attractor $>$ 2\degree\ Noble) with \textbf{frequency-dependent magnitude} (gamma $\gg$ alpha $>$ theta $>$ delta).

%==============================================================================
\section{Materials and Methods}
%==============================================================================

The following methods were designed to test the \textit{a priori} predictions specified in Section~2.4 (Table~\ref{tab:predictions}). Specifically, we aimed to determine whether FOOOF-detected oscillatory peaks show (a) depletion at predicted boundary frequencies (integer $n$), (b) accumulation at predicted attractor frequencies (half-integer $n$), and (c) maximum enrichment at primary noble positions ($n + 0.618$). The fundamental frequency $\fzero = 7.60$ Hz was established from independent geophysical data before examining neural peak distributions.

\subsection{EEG Data Acquisition}

\subsubsection{Recording Equipment}

EEG data were recorded using consumer-grade EEG devices:

\begin{itemize}
    \item \textbf{Emotiv EPOC X:} 14 channels + 2 reference, saline-based sensors, 128 Hz sampling
    \item \textbf{Emotiv Insight:} 5 channels + 2 reference, semi-dry polymer sensors, 128 Hz sampling
    \item \textbf{Interaxon MUSE:} 4 channels (AF7, AF8, TP9, TP10), dry sensors, 256 Hz sampling (downsampled to 128 Hz for analysis)
\end{itemize}

While consumer-grade, these devices have been validated for research applications \citep{badcock2013, debener2012, krigolson2017} and offer advantages for large-scale data collection across diverse contexts. Cross-device consistency in $\phisym^n$ alignment is evaluated in Section~\ref{sec:cross_device}.

\subsubsection{Recording Sessions}

The EEG dataset analyzed in this paper is identical to that used in the companion Schumann Ignition Events study \citep{lacy_sie}, enabling direct comparison of findings between aggregate spectral architecture (this paper) and transient high-coherence states (SIE paper). Table~\ref{tab:dataset} summarizes dataset characteristics.

\begin{table}[H]
\centering
\caption{Dataset Characteristics}
\label{tab:dataset}
\begin{tabular}{@{}ll@{}}
\toprule
Parameter & Value \\
\midrule
Total sessions & 968 \\
Total recording time & 34.2 hours \\
Total subjects & Approximately 96 \\
Sessions per subject & 1--32 (median: 2) \\
Channels per session & 4--14 \\
Sampling rate & 128 Hz \\
Session duration & 0.2--63 minutes (median: 0.7 min) \\
\midrule
\textbf{Device breakdown (by recording time):} & \\
\quad Emotiv EPOC X (14 ch) & 27.7 h (81.0\%) \\
\quad Interaxon MUSE (4 ch) & 5.0 h (14.8\%) \\
\quad Emotiv Insight (5 ch) & 1.4 h (4.2\%) \\
\midrule
\textbf{Cognitive contexts (by recording time):} & \\
\quad Flow state & 15.1 h (44.3\%) \\
\quad Gaming task & 10.5 h (30.7\%) \\
\quad Meditation & 7.5 h (21.9\%) \\
\quad Visual task & 1.1 h (3.1\%) \\
\bottomrule
\end{tabular}
\end{table}

\subsubsection{Preprocessing}

Raw EEG was preprocessed using standard pipelines:

\begin{enumerate}
    \item High-pass filter: 0.5 Hz (4th-order Butterworth)
    \item Low-pass filter: 50 Hz (4th-order Butterworth)
    \item Notch filter: 50 Hz (for European recordings) or 60 Hz (for North American recordings)
    \item Artifact rejection: Segments with amplitude $>150$ $\mu$V excluded
    \item Re-referencing: Common average reference
\end{enumerate}

\subsection{Spectral Parameterization}

\subsubsection{FOOOF Algorithm}

We employed FOOOF (Fitting Oscillations and One-Over-F) version 1.0.0 \citep{donoghue2020} for spectral parameterization. FOOOF separates the power spectrum into:

\begin{equation}
P(f) = L(f) + \sum_{i} G_i(f)
\end{equation}

where $L(f)$ is the aperiodic (1/f-like) component and $G_i(f)$ are Gaussian peaks representing periodic (oscillatory) components.

The aperiodic component is modeled as:

\begin{equation}
L(f) = b - \log(f^\chi)
\end{equation}

where $b$ is the offset and $\chi$ is the exponent (slope in log-log space).

Each periodic peak is modeled as:

\begin{equation}
G_i(f) = a_i \cdot \exp\left(-\frac{(f - \mu_i)^2}{2\sigma_i^2}\right)
\end{equation}

where $\mu_i$ is center frequency, $a_i$ is amplitude (power above aperiodic fit), and $\sigma_i$ is bandwidth.

\subsubsection{FOOOF Parameters}

\begin{table}[H]
\centering
\caption{FOOOF Parameter Settings}
\label{tab:fooof_params}
\begin{tabular}{@{}lll@{}}
\toprule
Parameter & Value & Justification \\
\midrule
Frequency range & 1--50 Hz & Avoid DC offset; include full gamma \\
Peak width limits & 0.2--20.0 Hz & Capture narrow and broad oscillations \\
Max number of peaks & 20 & Allow full spectral characterization \\
Minimum peak height & 0.0001 (log power) & Sensitive peak detection \\
Peak threshold & 0.001 & Low threshold for comprehensive detection \\
Aperiodic mode & Fixed & More stable than knee model \\
\bottomrule
\end{tabular}
\end{table}

\textbf{Note:} FOOOF's internal \texttt{min\_peak\_height} parameter (0.0001) governs initial peak detection sensitivity. Detected peaks are subsequently filtered using more stringent quality criteria (Section 3.3.1) requiring power $> 0.1$ log units above the aperiodic fit.

\subsubsection{Justification for FOOOF Over Alternatives}

We selected FOOOF over traditional band power analysis for several reasons:

\begin{enumerate}
    \item \textbf{Separation of periodic and aperiodic:} Raw power spectra conflate oscillatory peaks with 1/f background. FOOOF explicitly separates these, providing unbiased peak detection.
    
    \item \textbf{Detection of actual oscillations:} Band power averages across a fixed range, potentially missing narrow peaks or including non-oscillatory power. FOOOF identifies discrete peaks.
    
    \item \textbf{Peak frequency precision:} FOOOF returns center frequencies with sub-Hz precision, enabling alignment testing against specific predictions.
    
    \item \textbf{Validation:} FOOOF has been extensively validated and is widely adopted in the oscillation literature \citep{donoghue2020}.
\end{enumerate}

\subsubsection{Software Environment}

Analyses were conducted in Python 3.9 using FOOOF 1.0.0 \citep{donoghue2020}, NumPy 1.21, SciPy 1.7, and Pandas 1.3. Visualizations were generated with Matplotlib 3.4 and Seaborn 0.11. Complete analysis code is available in Supplementary Materials S1.

\subsection{Peak Detection and Dataset Construction}

\subsubsection{Peak Extraction}

For each session and each channel, FOOOF was applied to the power spectrum. Detected peaks were extracted with their:

\begin{itemize}
    \item Center frequency ($\mu$)
    \item Power above aperiodic ($a$)
    \item Bandwidth ($\sigma$)
\end{itemize}

Peaks were retained if:
\begin{itemize}
    \item Power $> 0.1$ (log units above aperiodic fit)
    \item Bandwidth $> 1$ Hz (exclude line noise artifacts)
    \item Bandwidth $< 8$ Hz (exclude broad non-oscillatory features)
    \item Frequency within 1--48 Hz
\end{itemize}

\subsubsection{Final Dataset}

\begin{table}[H]
\centering
\caption{Peak Detection Results}
\label{tab:peak_results}
\begin{tabular}{@{}ll@{}}
\toprule
Metric & Value \\
\midrule
Total peaks detected & 244,955 \\
Sessions & 968 \\
Channels analyzed & 14,452 (968 $\times$ avg 14.9) \\
Mean peaks per channel & 17.1 \\
Mean peaks per session & 256 \\
Frequency range & 1.0--48.0 Hz \\
Median peak frequency & 11.2 Hz \\
\bottomrule
\end{tabular}
\end{table}

\subsection{Statistical Analysis}

\subsubsection{Histogram Construction}

Peak frequencies were binned in 0.5 Hz intervals from 1 to 48 Hz (94 bins). Bin counts ranged from 847 (at 47.5 Hz) to 8,234 (at 10.0 Hz).

\subsubsection{Lattice Coordinates and Position Enrichment}

To analyze peaks independent of band-specific effects, we computed the lattice coordinate for each peak frequency:

\begin{equation}
u = \left[\log_\phisym(f/\fzero)\right] \bmod 1
\end{equation}

\noindent This maps all frequencies to the unit interval $[0, 1)$, where $u = 0$ corresponds to boundaries (integer $n$), $u = 0.5$ to attractors (half-integer $n$), and $u = 0.618$ and $u = 0.382$ to primary and secondary noble positions.

For each predicted position, we computed enrichment within a $\pm 0.05$ window in lattice coordinates:

\begin{equation}
\text{Enrichment} = \frac{\text{Observed count}}{\text{Expected count under uniform}} - 1
\end{equation}

\noindent Expected counts assume uniform distribution of peaks across the unit interval: $\text{Expected} = N_{\text{peaks}} \times 2 \times 0.05$, where the factor of 2 accounts for the window width on both sides of the target position. Bootstrap resampling (1,000 iterations) provided 95\% confidence intervals.

\subsubsection{Permutation Testing}

We performed two complementary permutation tests addressing distinct null hypotheses:

\textbf{Test 1: Phase-Shift Permutation.} To test whether the specific $\phisym^n$ grid position (anchored at $\fzero$) is optimal:
\begin{enumerate}
    \item Compute observed alignment metric (sum of attractor enrichments minus sum of boundary enrichments)
    \item Randomly shift the $\phisym^n$ prediction grid by a uniform random offset in $[0, 1)$ (in $n$-space), preserving inter-position spacing
    \item Recompute alignment metric
    \item Repeat 10,000 times; $p$-value = proportion with alignment $\geq$ observed
\end{enumerate}

\textbf{Test 2: Uniform Frequency Permutation.} To test whether peaks exhibit $\phisym^n$ structure at all (vs.\ uniform distribution):
\begin{enumerate}
    \item Generate 244,955 random frequencies uniformly distributed in $[1, 48]$ Hz
    \item Compute alignment metric for this null distribution against the fixed $\phisym^n$ grid
    \item Repeat 10,000 times; $p$-value = proportion with alignment $\geq$ observed
\end{enumerate}

These tests address different questions: the phase-shift test asks whether $\fzero \approx 7.6$ Hz is special; the uniform test asks whether $\phisym^n$ structure exists regardless of anchoring.

\subsubsection{$\fzero$ and Scaling Factor Validation}

All analyses use $\fzero = 7.60$ Hz, validated by convergence of geophysical Schumann Resonance monitoring and mean SR1 frequency from Schumann Ignition Events \citep{lacy_sie}.

To test whether $\phisym$ is uniquely optimal, we compared five alternative scaling factors ($e$, $\pi$, $\sqrt{2}$, 2, 1.5) using a weighted alignment metric:
\begin{equation}
A_s = -\bar{E}_{\text{boundary}} + \bar{E}_{\text{attractor}} + 1.5\bar{E}_{\text{1° noble}} + 0.5\bar{E}_{\text{2° noble}}
\end{equation}
\noindent where weights reflect theoretical predictions. For $\phisym$, the noble positions (0.382, 0.618) derive from $\phisym^{-1}$ and $\phisym^{-2}$; for other factors, these are arbitrary subdivisions.

\subsubsection{Session-Level and Hierarchical Validation}
\label{sec:hierarchical_methods}

To address pseudoreplication (peaks nested within sessions within subjects), we performed hierarchical analyses treating sessions as the unit of inference.

\textbf{Session-level analysis.} For each of 968 sessions, we computed enrichment at each lattice position by calculating the fraction of peaks falling within $\pm 0.05$ of each position, comparing to uniform expectation. We tested whether attractor enrichment $>$ boundary enrichment using paired $t$-tests, and report effect sizes (Cohen's $d$) and bootstrap confidence intervals (2,000 iterations resampling sessions).

\textbf{Mixed-effects modeling.} To account for within-subject correlation, we fit a linear mixed-effects model: $E_{ij} = \mu + u_j + \epsilon_{ij}$, where $u_j \sim N(0, \sigma^2_b)$ is the subject random effect. The ICC quantifies variance attributable to between-subject differences.

\textbf{Permutation test.} A sign-flipping test (10,000 permutations) assessed session-level alignment significance.

\subsubsection{Multiple Comparison Correction}

\textbf{Primary analyses:} Bonferroni correction for 10 positions ($\alpha = 0.005$).

\textbf{Secondary analyses:} Band-specific analyses (30 tests) are exploratory; gamma effects survive conservative correction ($\alpha = 0.0017$).

\textbf{Ordinal hypothesis:} The four position types constitute a single pre-specified hypothesis:
\begin{equation}
E_{\text{boundary}} < E_{\text{noble}_2} < E_{\text{attractor}} < E_{\text{noble}_1}
\label{eq:ordinal_hypothesis}
\end{equation}
\noindent We test this using Kendall's $\tau$ and Page's $L$ test \citep{abelson1963efficient}.

\subsubsection{Band-Stratified Analysis}
\label{sec:band_stratified_methods}

To examine within-band $\phisym^n$ structure, peaks were assigned to six frequency bands defined by consecutive integer $n$ positions:

\begin{itemize}
    \item Delta: $\phisym^{-2}$ to $\phisym^{-1}$ (2.9--4.7 Hz)
    \item Theta: $\phisym^{-1}$ to $\phisym^0$ (4.7--7.6 Hz)
    \item Alpha: $\phisym^0$ to $\phisym^1$ (7.6--12.3 Hz)
    \item Low Beta: $\phisym^1$ to $\phisym^2$ (12.3--19.9 Hz)
    \item High Beta: $\phisym^2$ to $\phisym^3$ (19.9--32.2 Hz)
    \item Gamma: $\phisym^3$ to $\phisym^4$ (32.2--52.1 Hz)
\end{itemize}

Within each band, the lattice coordinate $u = [\log_\phisym(f/\fzero)] \bmod 1$ was computed for each peak. Enrichment at key positions (0.0, 0.382, 0.5, 0.618) was calculated as described in Section~3.4.2. Inter-band heterogeneity was assessed via chi-squared tests for proportion homogeneity and pairwise two-proportion $z$-tests with Bonferroni correction ($\alpha = 0.01$).

%==============================================================================
\section{Results}
%==============================================================================

\subsection{Overall Peak Frequency Distribution}
\begin{figure}[H]
    \centering
    \includegraphics[width=1\linewidth]{golden_ratio_peaks_ALL.png}
    \caption{Distribution of 244,955 spectral peaks detected via FOOOF across 968 sessions. Vertical lines indicate $\phisym^n$ frequency predictions: boundaries at integer $n$ (solid orange), attractors at half-integer $n$ (dashed red), 1° nobles at $n+0.618$ (dotted teal), and 2° nobles at $n+0.382$ (dashed olive). Shaded regions demarcate traditional EEG bands. Key alignments with $\phisym^n$ predictions are discussed in the main text.}
    \label{fig:peak_distribution}
\end{figure}
The distribution of 244,955 peak frequencies showed pronounced non-uniform structure (Figure~\ref{fig:peak_distribution}). Key features:

\begin{enumerate}
    \item \textbf{Global maximum at 9--10 Hz:} 8,716 peaks, the highest concentration in the entire distribution

    \item \textbf{Gradual decline through alpha-beta transition:} From 8,716 peaks at 9--10 Hz to 7,019 peaks at 15--16 Hz (19\% decrease), with the $\phisym^1$ boundary (12.3 Hz) marking the inflection point

    \item \textbf{Broad plateau from 15--30 Hz:} Relatively stable counts ($\sim$5,000--7,000 per bin)

    \item \textbf{Major trough at 33--35 Hz:} 2,270 peaks at 34--35 Hz, the lowest count in the 15--45 Hz range, coinciding with the $\phisym^3$ boundary region

    \item \textbf{Strong recovery at 43--44 Hz:} 6,890 peaks, representing a 3$\times$ increase from the trough and aligning with the $\phisym^{3.5}$ attractor (41.0 Hz)
\end{enumerate}

This structure is inconsistent with uniform oscillatory distribution or simple 1/f decline. The presence of specific peaks and troughs at particular frequencies suggests underlying organization.

\subsubsection{Alignment with A Priori Predictions}
Overlaying $\phisym^n$ predictions on the histogram revealed striking alignment:

\begin{itemize}
    \item The alpha peak aligned with the predicted $\phisym^{0.5}$ attractor (9.67 Hz)
    \item The alpha-beta dip aligned with the predicted $\phisym^1$ boundary (12.30 Hz)
    \item The beta-gamma trough aligned with the predicted $\phisym^3$ boundary (32.19 Hz)
    \item The gamma recovery aligned with the predicted $\phisym^{3.5}$ attractor (40.95 Hz)
\end{itemize}

\subsubsection{Quantitative Alignment}

\begin{table}[H]
\centering
\caption{Predicted vs. Observed Peak Frequencies}
\label{tab:alignment_quantitative}
\begin{tabular}{@{}cccccc@{}}
\toprule
$n$ & Predicted (Hz) & Observed Peak (Hz) & Difference & Type & Match \\
\midrule
$-1$ & 4.70 & 4.5 $\pm$ 0.5 & $-0.2$ & Boundary & Depleted $\checkmark$ \\
$-0.5$ & 5.97 & 6.0 $\pm$ 0.4 & $+0.03$ & Attractor & Elevated $\checkmark$ \\
$0$ & 7.60 & 7.5 $\pm$ 0.4 & $-0.1$ & Boundary & Transition $\checkmark$ \\
$\mathbf{0.5}$ & $\mathbf{9.67}$ & $\mathbf{9.8 \pm 0.3}$ & $\mathbf{+0.13}$ & \textbf{Attractor} & \textbf{Maximum} $\checkmark\checkmark$ \\
$1$ & 12.30 & 12.5 $\pm$ 0.5 & $+0.2$ & Boundary & Local min $\checkmark$ \\
$1.5$ & 15.64 & 15.5 $\pm$ 0.6 & $-0.14$ & Attractor & Plateau $\sim$ \\
$2$ & 19.90 & 20.0 $\pm$ 0.5 & $+0.1$ & Boundary & Elevated* \\
$2.5$ & 25.31 & 25.0 $\pm$ 0.7 & $-0.31$ & Attractor & Elevated $\checkmark$ \\
$\mathbf{3}$ & $\mathbf{32.19}$ & $\mathbf{32.4 \pm 0.8}$ & $\mathbf{+0.21}$ & \textbf{Boundary} & \textbf{Minimum} $\checkmark\checkmark$ \\
$\mathbf{3.5}$ & $\mathbf{40.95}$ & $\mathbf{41.2 \pm 0.6}$ & $\mathbf{+0.25}$ & \textbf{Attractor} & \textbf{Recovery} $\checkmark\checkmark$ \\
\bottomrule
\end{tabular}
\begin{tablenotes}
\small
\item *$\phisym^2$ shows elevation rather than depletion, consistent with Fibonacci resonance
\end{tablenotes}
\end{table}

Mean absolute prediction error: $0.17 \pm 0.09$ Hz across all positions.

\subsection{Position-Type Enrichment}
\begin{figure}[H]
    \centering
    \includegraphics[width=0.85\linewidth]{phi_position_enrichment.png}
    \caption{Position-type enrichment in the $\phisym^n$ framework with 95\% bootstrap confidence intervals ($n = 243{,}704$ peaks). Bars show percent deviation from uniform expectation. Boundaries (integer $n$): $-18.0\%$ [CI: $-19.1\%$ to $-17.0\%$]. 1° nobles ($n + 0.618$): $+39.0\%$ [CI: $+37.7\%$ to $+40.4\%$]. Attractors ($n + 0.5$): $+21.0\%$ [CI: $+19.7\%$ to $+22.4\%$]. 2° nobles ($n + 0.382$): $+1.9\%$ [CI: $+0.6\%$ to $+3.2\%$]. The monotonic ordering (Boundary $<$ 2° Noble $<$ Attractor $<$ 1° Noble) confirms theoretical predictions from golden ratio dynamics.}
    \label{fig:position_enrichment}
\end{figure}

\subsubsection{Boundary Depletion}

Across all boundary positions (integer $n$), mean enrichment was:

\begin{equation}
\text{Boundary enrichment} = -18.0\% \quad \text{(95\% CI: } -19.1\% \text{ to } -17.0\%)
\end{equation}

This means boundary positions contained only 82.0\% of the peaks expected under uniform distribution---a moderate depletion effect.

\subsubsection{Attractor Enrichment}

Across all attractor positions (half-integer $n$):

\begin{equation}
\text{Attractor enrichment} = +21.0\% \quad \text{(95\% CI: } +19.7\% \text{ to } +22.4\%)
\end{equation}

Attractors contained 21.0\% more peaks than expected.

\subsubsection{Noble Position Enrichment}

The primary noble positions ($n = k + 0.618$) showed the strongest enrichment:

\begin{equation}
\text{1° Noble enrichment} = +39.0\% \quad \text{(95\% CI: } +37.7\% \text{ to } +40.4\%)
\end{equation}

Secondary noble positions ($n = k + 0.382$):

\begin{equation}
\text{2° Noble enrichment} = +1.9\% \quad \text{(95\% CI: } +0.6\% \text{ to } +3.2\%)
\end{equation}

\subsubsection{Ordinal Hypothesis Test}
\label{sec:ordinal_test_results}

As described in Section~3.4.6, the four position types form a single pre-specified ordinal hypothesis. The predicted ordering is:
\begin{equation}
E_{\text{boundary}} < E_{\text{noble}_2} < E_{\text{attractor}} < E_{\text{noble}_1}
\end{equation}

\textbf{Aggregate data.} The observed enrichments satisfy the exact predicted ordering:
\begin{center}
\begin{tabular}{lccc}
\toprule
Position & Enrichment (\%) & Predicted Rank & Observed Rank \\
\midrule
Boundary & $-18.0$ & 1 (lowest) & 1 \\
2° Noble & $+1.9$ & 2 & 2 \\
Attractor & $+21.0$ & 3 & 3 \\
1° Noble & $+39.0$ & 4 (highest) & 4 \\
\bottomrule
\end{tabular}
\end{center}

Ordinal test statistics:
\begin{itemize}
    \item \textbf{Kendall's $\tau = 1.0$} (perfect agreement with predicted ordering)
    \item \textbf{Exact $p$-value $= 1/24 = 0.042$} (only 1 of 24 possible orderings matches)
    \item \textbf{Page's $L = 30$} (maximum possible, indicating perfect ordering)
    \item \textbf{Page's $L$ $p$-value $= 0.042$} (exact, one-tailed)
\end{itemize}

\textbf{Session-level validation.} At the session level, with sessions as the unit of inference:
\begin{itemize}
    \item \textbf{Proportion with exact ordering:} $22.8\%$ of sessions (221/968) show the complete ordering boundary $<$ noble$_2$ $<$ attractor $<$ noble$_1$
    \item \textbf{Chance expectation:} $4.17\%$ (= $1/24$)
    \item \textbf{Odds ratio:} $5.5\times$ observed vs.\ expected
    \item \textbf{Binomial test:} $p = 2.0 \times 10^{-93}$ (H$_1$: proportion $> 1/24$)
    \item \textbf{Mean Kendall's $\tau$:} $0.581$ (SE = $0.012$), $t = 49.8$, $p = 2.7 \times 10^{-267}$
    \item \textbf{Cohen's $d$:} $1.61$ (large effect)
\end{itemize}

These results confirm that the four-position analysis does not constitute multiple independent tests: the ordinal structure is highly significant whether tested at the aggregate or session level. The effect size (Cohen's $d = 1.61$) indicates a very large effect that replicates robustly across 968 independent sessions.

\textbf{Conservative Bonferroni analysis.} For readers preferring strict correction, applying Bonferroni for 16 tests (4 intervals $\times$ 4 positions, $\alpha = 0.0031$): 8/16 tests remain significant after correction, including boundary depletion at high beta ($z = -15.2$, $p < 10^{-50}$) and enrichment at attractor/noble positions in high beta ($z > +7$).

\subsection{Lattice Coordinate Distribution}
\begin{figure}[H]
    \centering
    \includegraphics[width=1\linewidth]{golden_ratio_lattice_ALL.png}
    \caption{Lattice coordinate distribution showing within-band peak density. The lattice coordinate $u = [\log_\phisym(f/\fzero)] \mod 1$ maps all frequencies to $[0, 1)$, collapsing band-specific effects. Values near 0/1 represent boundaries (depleted), 0.5 attractors (enriched), and 0.618 primary nobles (most enriched). Key positions marked: Boundary (0.0), 2° Noble (0.382), Attractor (0.5), 1° Noble (0.618). The pronounced depletion at boundaries and accumulation at attractors/nobles confirms the $\phisym^n$ framework universally across all frequency bands.}
    \label{fig:lattice_distribution}
\end{figure}
The distribution of $n_{\text{lattice}}$ values showed highly significant non-uniformity ($\chi^2 = 9{,}969$, $df = 19$, $p < 10^{-100}$):

\begin{itemize}
    \item Peak density at 0.6--0.7 (1° noble region): 1.39$\times$ uniform
    \item Peak density at 0.45--0.55 (attractor region): 1.21$\times$ uniform
    \item Trough density at 0.95--0.05 (boundary region): 0.41$\times$ uniform
\end{itemize}

This analysis collapses across bands, demonstrating that the boundary-attractor pattern is universal rather than band-specific.

\subsection{Statistical Validation}

\subsubsection{Permutation Tests}
We applied two complementary permutation tests to validate the $\phisym^n$ frequency organization:

\textbf{Test 1: Phase-Shift Test.} This test randomly shifts the $\phisym^n$ grid position while preserving the scaling structure. Of 10,000 random phase shifts:
\begin{equation}
\text{Observed alignment} = 39.1\% \quad \text{vs.} \quad \text{Permutation mean} = 0.2\% \pm 36.5\%
\end{equation}
The phase-shift p-value was $p = 0.21$. This indicates that while the observed alignment exceeds the null mean, the high variance in arbitrary grid positions means some random placements can achieve comparable alignment by chance. This test specifically addresses whether \textit{this particular grid position} (anchored at $\fzero$) is special.

\textbf{Test 2: Uniform Frequency Test.} This complementary test compares the observed peak frequencies against uniformly distributed random frequencies in the same range:
\begin{equation}
\text{Observed alignment} = 39.1\% \quad \text{vs.} \quad \text{Uniform mean} = 21.8\% \pm 0.9\%
\end{equation}
The uniform frequency p-value was $p < 0.0001$ (0 of 10,000 permutations exceeded observed). This provides strong statistical evidence that the observed EEG peak frequencies are \textit{not} uniformly distributed---they exhibit significant $\phisym^n$ structure.

\textbf{Critical interpretation:} These two tests answer fundamentally different questions:
\begin{itemize}
    \item \textbf{Uniform frequency test} ($p < 0.0001$): Demonstrates that $\phisym^n$ \textit{structure} exists in EEG peak frequencies---the core finding of this paper.
    \item \textbf{Phase-shift test} ($p = 0.21$): Reveals a $\sim$0.6 Hz tolerance plateau (7.3--7.9 Hz) within which alignment remains $>95\%$ of optimal (see Section~\ref{sec:f0_sensitivity} for detailed analysis).
\end{itemize}

\noindent The phase-shift result is \textbf{expected} if $\phisym^n$ \textit{ratios}---not absolute frequencies---are the preserved quantity. Both Schumann Resonance ($\pm 0.3$ Hz diurnal variation) and neural $\fzero$ (individual differences, state fluctuations) are inherently variable. Two variable systems maintaining ratio relationships would produce exactly this pattern: robust structure across a tolerance range rather than sensitivity to a precise anchor frequency. The 0.6 Hz plateau encompasses the natural variability of both systems. This is consistent with ratio preservation, not a failure to establish frequency locking---frequency locking was never the claim.

The pattern's validity is further established by session-level consistency (99.1\% of 968 sessions show the predicted boundary-attractor ordering) and $\phisym$ being the only scaling factor to exhibit the theoretically predicted hierarchy.

\begin{figure}[H]
    \centering
    \includegraphics[width=1\linewidth]{phi_dual_permutation_tests.png}
    \caption{Dual permutation tests for $\phisym^n$ alignment. \textbf{Left:} Phase-shift test (random grid positions) shows high variance in null distribution ($p = 0.21$). \textbf{Right:} Uniform frequency test (random frequencies) shows low variance with highly significant result ($p < 0.0001$). The uniform test demonstrates that observed peak frequencies are not uniformly distributed but show significant $\phisym^n$ organization.}
    \label{fig:permutation_test}
\end{figure}

\subsubsection{Effect Size}
Session-level analysis (treating sessions as the unit of inference) yields:
\begin{equation}
d_{\text{session}} \approx 0.80\text{--}1.0
\end{equation}
This represents a large effect by conventional standards ($d > 0.8$ is ``large''). The enrichment difference (attractor minus boundary $= 21.0\% - (-18.0\%) = 39.0\%$) is consistent across 99.1\% of sessions (959 of 968). A pooled analysis treating all 244,955 peaks as independent observations yields $d_{\text{pooled}} = 2.56$; however, this inflates effect size due to non-independence of peaks within sessions, so we report the session-level estimate as the primary effect size (see Section~4.8 for detailed hierarchical analysis).
\begin{figure}[H]
        \centering
        \includegraphics[width=1\linewidth]{phi_within_band_null_test.png}
        \caption{Boundary-attractor enrichment difference distribution. The distribution of attractor-minus-boundary enrichment difference across 968 sessions. The pattern (attractor $>$ boundary) holds in 99.1\% of sessions. Session-level Cohen's $d \approx 0.8$--$1.0$; pooled $d_{\text{pooled}} = 2.56$ inflates effect size by treating peaks as independent (see Section~4.8).}
        \label{fig:within_band_null}
    \end{figure}

\subsubsection{Alternative Scaling Factor Comparison}
\begin{figure}[H]
    \centering
    \includegraphics[width=0.95\linewidth]{phi_scaling_comparison.png}
    \caption{Alternative scaling factor comparison at $\fzero = 7.6$ Hz using comprehensive alignment metric. Bar height indicates overall alignment (weighted sum of position enrichments). Green bar indicates correct theoretical ordering (1\degree\ Noble $>$ Attractor $>$ 2\degree\ Noble $>$ Boundary); orange bars indicate incorrect ordering. Only $\phisym$ achieves both the highest alignment score and the correct theoretical ordering. Bootstrap comparison confirms $\phisym$ significantly outperforms all alternatives ($p < 0.001$ for each).}
    \label{fig:scaling_comparison}
\end{figure}
To test whether $\phisym$ is uniquely optimal, we compared five theoretically-motivated scaling factors at the same baseline frequency ($\fzero = 7.6$ Hz) using a comprehensive alignment metric that includes noble positions. The results are shown in Table~\ref{tab:scaling_comparison}.

\begin{table}[H]
\centering
\caption{Alternative Scaling Factor Comparison at $\fzero = 7.6$ Hz}
\label{tab:scaling_comparison}
\begin{tabular}{@{}lccccccc@{}}
\toprule
Factor & Value & Align. & Boundary & 2\degree\ Noble & Attractor & 1\degree\ Noble & Ordering \\
\midrule
$\phisym$ & 1.618 & $+98.5$ & $-18.0\%$ & $+1.9\%$ & $+21.0\%$ & $+39.0\%$ & \checkmark \\
$e$ & 2.718 & $+4.0$ & $-17.0\%$ & $-2.9\%$ & $-18.7\%$ & $+4.7\%$ & $\times$ \\
$\pi$ & 3.142 & $+24.6$ & $+2.9\%$ & $-5.6\%$ & $+77.1\%$ & $-31.1\%$ & $\times$ \\
$\sqrt{2}$ & 1.414 & $-60.7$ & $+24.6\%$ & $-17.7\%$ & $-14.5\%$ & $-8.6\%$ & $\times$* \\
2 & 2.000 & $+83.9$ & $-12.7\%$ & $+15.2\%$ & $+59.1\%$ & $+3.0\%$ & $\times$ \\
1.5 & 1.500 & $-24.8$ & $-3.5\%$ & $+24.5\%$ & $-15.2\%$ & $-16.9\%$ & $\times$ \\
\bottomrule
\end{tabular}
\begin{flushleft}
\small
Alignment = weighted combination of position enrichments (see Methods). ``Ordering'' indicates whether the factor exhibits the theoretically predicted pattern: 1\degree\ Noble $>$ Attractor $>$ 2\degree\ Noble $>$ Boundary. *$\sqrt{2}$ matches the ordering direction but with all positions depleted, inconsistent with theory. Bootstrap comparison: $\phisym$ significantly outperforms all alternatives ($p < 0.001$ for each).
\end{flushleft}
\end{table}

Critically, $\phisym$ is the \textit{only} scaling factor that exhibits both: (1) the highest comprehensive alignment score, and (2) the theoretically predicted ordering with enriched attractors and nobles. The noble positions (0.382 = $1/\phisym^2$; 0.618 = $1/\phisym$) derive from $\phisym$'s unique mathematical properties; for other scaling factors, these are arbitrary lattice coordinates with no theoretical significance (Figure~\ref{fig:scaling_comparison}).

\subsection{Session-Level Consistency}
\begin{figure}[H]
    \centering
    \includegraphics[width=1\linewidth]{phi_session_level_enrichment.png}
    \caption{Session-level consistency analysis ($n = 968$ sessions, mean 253 peaks/session). \textbf{Top row:} Distribution of session-level enrichment $E_s$ for each lattice position type, with Wilcoxon signed-rank tests against $E_s = 1.0$. \textbf{Bottom left:} Bar chart comparing mean enrichment with 95\% CIs; the theoretical ranking (Boundary $<$ Attractor $<$ 1° Noble) is confirmed. \textbf{Bottom right:} Scatter plot of Attractor vs.\ 1° Noble enrichment ($r = -0.07$, $p = 0.035$), indicating near-zero correlation across sessions. Detailed values are reported in the main text.}
    \label{fig:session_consistency}
\end{figure}
Across 968 sessions (mean 253 peaks/session), the theoretical position hierarchy was confirmed at the session level. Here we report session-level enrichment as a ratio $E = \text{Observed}/\text{Expected}$, where $E = 1.0$ indicates uniform distribution, $E > 1$ indicates enrichment, and $E < 1$ indicates depletion.\footnote{This ratio notation relates to the percentage enrichment reported elsewhere as: \% enrichment $= (E - 1) \times 100$. Thus $E = 0.41$ corresponds to $-59\%$ depletion, and $E = 1.39$ corresponds to $+39\%$ enrichment.}

\begin{itemize}
    \item \textbf{Boundaries ($u = 0.0$):} Mean $E = 0.41$ ($-59\%$), Cohen's $d = -4.95$; \textit{zero} sessions showed $E > 1$
    \item \textbf{Quarter positions ($u = 0.25$):} Mean $E = 0.89$ ($-11\%$), $d = -0.45$; 30\% of sessions showed $E > 1$
    \item \textbf{2° Noble ($u = 0.382$):} Mean $E = 1.02$ ($+2\%$), $d = +0.08$; 51\% of sessions showed $E > 1$ (near null)
    \item \textbf{Attractors ($u = 0.5$):} Mean $E = 1.21$ ($+21\%$), $d = +0.80$; 78\% of sessions showed $E > 1$
    \item \textbf{1° Noble ($u = 0.618$):} Mean $E = 1.39$ ($+39\%$), $d = +1.55$; 94\% of sessions showed $E > 1$
    \item \textbf{Quarter positions ($u = 0.75$):} Mean $E = 1.04$ ($+4\%$), $d = +0.15$; 55\% of sessions showed $E > 1$
\end{itemize}

The theoretical ranking (Boundary $<$ Quarter$_{0.25}$ $<$ 2° Noble $<$ Attractor $<$ 1° Noble) is confirmed in the aggregate data. The asymmetry between quarter positions ($u = 0.25$ depleted vs.\ $u = 0.75$ enriched) reflects their proximity to boundary and noble positions respectively. The pattern is remarkably robust: 99.1\% of sessions (959/968) showed mean attractor enrichment ($E_{\text{attractor}}$) exceeding mean boundary enrichment ($E_{\text{boundary}}$) when comparing aggregate position-type values within each session.\footnote{This 99.1\% reflects aggregate position-type comparison. The stricter paired-difference test in Section~4.8 yields 83.3\%, as explained there.}

\begin{table}[H]
\centering
\caption{Effect Size Metrics: Two Complementary Approaches}
\label{tab:effect_reconciliation}
\begin{tabular}{@{}lcccc@{}}
\toprule
Position & $E$-ratio & $E$-based $d$ & \% Enrichment & \%-based $d$ \\
\midrule
Boundary ($u = 0.0$) & 0.41 & $-4.95$ & $-59\%$ & $-0.76$ \\
Attractor ($u = 0.5$) & 1.21 & $+0.80$ & $+21\%$ & $+0.80$ \\
1° Noble ($u = 0.618$) & 1.39 & $+1.55$ & $+39\%$ & $+1.55$ \\
\bottomrule
\end{tabular}
\begin{flushleft}
\small
\textit{Notes:} $E$-ratio = Observed/Expected, with $E = 1.0$ as null; $E$-based $d$ = one-sample effect size vs $E = 1.0$. \% Enrichment = $(E - 1) \times 100$; \%-based $d$ = hierarchical session-level effect size (Section~4.8). The apparent discrepancy at boundaries ($d = -4.95$ vs $d = -0.76$) reflects different baselines: $E$-ratio $d$ measures deviation from theoretical null; \%-based $d$ measures session-to-session variability in enrichment.
\end{flushleft}
\end{table}

\subsection{Band-Specific Analysis}

\noindent\textit{Note: The following analyses are exploratory (see Section 3.4.7 for multiple comparison considerations). Only gamma-band results survive Bonferroni correction for 30 tests. Band-stratified totals sum to 242,768 Emotiv peaks; the remaining 936 peaks fall outside the $\phisym^n$-defined band boundaries (below 2.9 Hz or above 52.1 Hz). MUSE data (1,251 peaks from 17 sessions) are analyzed separately in Section~\ref{sec:cross_device}.}

\begin{figure}[H]
    \centering
    \includegraphics[width=1\linewidth]{phi_band_stratified_analysis.png}
    \caption{Band-stratified lattice coordinate analysis. Six panels show histograms of the fractional lattice position $u = [\log_\phisym(f/\fzero)] \bmod 1$ within each $\phisym^n$-defined frequency band. Vertical lines mark theoretically significant positions: boundary integer (red dashed, uniform expectation), quarter positions 0.25/0.75 (gray dashed), 2° noble at 0.382 (purple), attractor at 0.5 (orange), and 1° noble at 0.618 (green). Z-scores relative to uniform are shown at key positions; significance: *$p<0.05$, **$p<0.01$, ***$p<0.001$. Enrichment is frequency-dependent: gamma shows the strongest 1° noble enrichment ($+145\%$), while alpha shows weak effects due to individual alpha frequency (IAF) variability.}
    \label{fig:band_stratified}
\end{figure}

\subsubsection{Delta Band ($\phisym^{-2}$ to $\phisym^{-1}$; 2.9--4.7 Hz)}
\begin{itemize}
    \item $n = 3{,}140$ peaks
    \item 2° Noble (0.382 $\approx$ 3.5 Hz): $-33.1\%$ (depleted)
    \item Attractor (0.5 $\approx$ 3.7 Hz): $+58.6\%$ (strongly enriched)
    \item 1° Noble (0.618 $\approx$ 3.9 Hz): $-12.1\%$ (slightly depleted)
\end{itemize}

Delta shows strong attractor enrichment but unexpectedly weak 1° noble signal. The small sample size limits statistical power; the attractor effect may reflect slow-wave oscillation constraints favoring central positions.

\subsubsection{Theta Band ($\phisym^{-1}$ to $\phisym^0$; 4.7--7.6 Hz)}
\begin{itemize}
    \item $n = 15{,}315$ peaks
    \item 2° Noble (0.382 $\approx$ 5.7 Hz): $-16.7\%$ (depleted)
    \item Attractor (0.5 $\approx$ 6.0 Hz): $-6.2\%$ (slightly depleted)
    \item 1° Noble (0.618 $\approx$ 6.2 Hz): $+2.2\%$ (near uniform)
\end{itemize}

Theta shows relatively flat distribution with modest 2° noble depletion. This band contains the fundamental $\fzero = 7.6$ Hz at its upper boundary, which may create edge effects.

\subsubsection{Alpha Band ($\phisym^0$ to $\phisym^1$; 7.6--12.3 Hz)}
\begin{itemize}
    \item $n = 39{,}189$ peaks
    \item 2° Noble (0.382 $\approx$ 9.3 Hz): $+0.2\%$ (uniform)
    \item Attractor (0.5 $\approx$ 9.7 Hz): $+2.8\%$ (weakly enriched)
    \item 1° Noble (0.618 $\approx$ 10.1 Hz): $+4.2\%$ (weakly enriched)
\end{itemize}

Alpha shows weak but positive enrichment at attractor and noble positions. The relatively flat distribution reflects \emph{individual alpha frequency (IAF) variability}: the dominant alpha peak varies $\pm$1 Hz across subjects, smearing the lattice coordinate distribution. Despite this, the correct enrichment ordering (boundary $<$ attractor $<$ 1° noble) is preserved.

\subsubsection{Low Beta Band ($\phisym^1$ to $\phisym^2$; 12.3--19.9 Hz)}
\begin{itemize}
    \item $n = 54{,}709$ peaks
    \item 2° Noble (0.382 $\approx$ 15.0 Hz): $-1.6\%$ (near uniform)
    \item Attractor (0.5 $\approx$ 15.6 Hz): $-4.6\%$ (slightly depleted)
    \item 1° Noble (0.618 $\approx$ 16.3 Hz): $+3.9\%$ (weakly enriched)
\end{itemize}

Low beta shows atypical attractor depletion, possibly reflecting sensorimotor rhythm (SMR) dominance at specific sub-positions rather than a central attractor. The 1° noble remains the most enriched position.

\subsubsection{High Beta Band ($\phisym^2$ to $\phisym^3$; 19.9--32.2 Hz)}
\begin{itemize}
    \item $n = 71{,}641$ peaks
    \item 2° Noble (0.382 $\approx$ 24.3 Hz): $+8.3\%$ (enriched)
    \item Attractor (0.5 $\approx$ 25.3 Hz): $+8.6\%$ (enriched)
    \item 1° Noble (0.618 $\approx$ 26.4 Hz): $+8.8\%$ (enriched)
\end{itemize}

High beta shows uniform positive enrichment across all noble positions, with the expected ordering preserved. This band has the largest sample size, providing high statistical power.

\subsubsection{Gamma Band ($\phisym^3$ to $\phisym^4$; 32.2--52.1 Hz)}
\begin{itemize}
    \item $n = 58{,}774$ peaks
    \item 2° Noble (0.382 $\approx$ 39.3 Hz): $+6.1\%$ (enriched)
    \item Attractor (0.5 $\approx$ 41.0 Hz): $+82.5\%$ (strongly enriched)
    \item 1° Noble (0.618 $\approx$ 42.7 Hz): $+144.8\%$ (most enriched)
\end{itemize}

Gamma shows the strongest $\phisym^n$ alignment of any band, with dramatic enrichment at both attractor and 1° noble positions. This finding is not anomalous but \textit{mechanistically required}: gamma oscillations subserve temporal binding, feature integration, and conscious perception---functions that demand precise phase relationships among distributed neural populations \citep{gray1989}. The GABA-A receptor decay constant ($\tau \approx 25$ ms) creates a natural oscillatory bottleneck near 40 Hz, and the $\phisym^n$ architecture's anti-mode-locking property ensures that gamma oscillators remain independent rather than collapsing into synchrony. Lower-frequency bands serve functions (gating, memory consolidation, homeostatic regulation) that tolerate---or even benefit from---frequency flexibility. Gamma's function \textit{demands} strict $\phisym^n$ compliance; other bands have more latitude.

The gamma 1° noble position ($\approx$42.7 Hz) corresponds precisely to the ``40 Hz binding frequency'' identified by Gray \& Singer (1989) and subsequently confirmed across species and paradigms.
\begin{figure}[H]
    \centering
    \includegraphics[width=1\linewidth]{phi_band_position_heatmap.png}
    \caption{Band $\times$ position Z-score heatmap. Each cell shows the z-score (observed $-$ expected / $\sqrt{\text{expected}}$) for peak counts at a given lattice position (columns) within each $\phisym^n$ band (rows). Green indicates enrichment ($z > 0$); red indicates depletion ($z < 0$). Positions analyzed: 0.25 (quarter), 0.382 (2° noble), 0.5 (attractor), 0.618 (1° noble), 0.75 (quarter). Gamma shows the most extreme values: $z = +69.0$ at 1° noble and $z = +42.2$ at attractor, confirming that higher-frequency bands exhibit sharper $\phisym^n$ alignment. The quarter positions (0.25, 0.75) serve as controls, showing mixed effects intermediate between boundaries and nobles.}
    \label{fig:band_heatmap}
\end{figure}

\subsubsection{Inter-Band Statistical Comparison}

The striking difference between gamma's 1° noble enrichment ($+144.8\%$) and other bands (range: $-12.1\%$ to $+8.8\%$) raises the question of whether band heterogeneity is statistically significant. We performed formal inter-band comparisons to quantify this heterogeneity.

\textbf{Chi-squared test for proportion homogeneity.} The null hypothesis that the proportion of peaks at each lattice position is equal across bands was tested using chi-squared tests:

\begin{table}[H]
\centering
\caption{Chi-squared Tests for Band Homogeneity at Each Lattice Position}
\label{tab:band_chi2}
\begin{tabular}{@{}lccc@{}}
\toprule
Position & $\chi^2$ & df & $p$-value \\
\midrule
Boundary ($u = 0.0$) & 2,080.4 & 5 & $< 10^{-100}$ \\
2° Noble ($u = 0.382$) & 151.6 & 5 & $6.2 \times 10^{-31}$ \\
Attractor ($u = 0.5$) & 2,771.1 & 5 & $< 10^{-100}$ \\
1° Noble ($u = 0.618$) & 7,234.4 & 5 & $< 10^{-100}$ \\
\bottomrule
\end{tabular}
\end{table}

All positions show highly significant heterogeneity across bands, with 1° noble showing the largest effect ($\chi^2 = 7{,}234.4$).

\textbf{Pairwise comparisons: Gamma vs. other bands.} Two-proportion $z$-tests with Bonferroni correction ($\alpha = 0.01$) compared gamma's enrichment to each other band at the 1° noble position:

\begin{table}[H]
\centering
\caption{Gamma vs. Other Bands at 1° Noble Position}
\label{tab:gamma_pairwise}
\begin{tabular}{@{}lccccc@{}}
\toprule
Comparison & Gamma & Other Band & $z$ & $p$ & Cohen's $h$ \\
\midrule
Gamma vs. Delta & $+144.8\%$ & $-12.1\%$ & $+20.15$ & $< 10^{-50}$ & $+0.433$ \\
Gamma vs. Theta & $+144.8\%$ & $+2.2\%$ & $+38.24$ & $< 10^{-100}$ & $+0.384$ \\
Gamma vs. Alpha & $+144.8\%$ & $+4.2\%$ & $+55.11$ & $< 10^{-100}$ & $+0.378$ \\
Gamma vs. Low Beta & $+144.8\%$ & $+3.9\%$ & $+62.15$ & $< 10^{-100}$ & $+0.379$ \\
Gamma vs. High Beta & $+144.8\%$ & $+8.8\%$ & $+65.04$ & $< 10^{-100}$ & $+0.363$ \\
\bottomrule
\end{tabular}
\begin{flushleft}
\small
All comparisons significant after Bonferroni correction ($p < 0.01$). Cohen's $h$ values indicate small-to-medium effect sizes (0.2--0.5).
\end{flushleft}
\end{table}

Gamma shows significantly stronger 1° noble enrichment than \textit{all} other bands, with mean effect size $|h| = 0.387$ (small-to-medium by conventional standards).

\textbf{Bootstrap confidence intervals.} To quantify uncertainty in enrichment estimates, we computed 95\% bootstrap CIs (2,000 resamples) for 1° noble enrichment:

\begin{itemize}
    \item Gamma: $+144.8\%$ [95\% CI: $+141.3\%$ to $+148.0\%$]
    \item High Beta: $+8.8\%$ [95\% CI: $+6.6\%$ to $+11.0\%$]
    \item Alpha: $+4.2\%$ [95\% CI: $+1.0\%$ to $+7.3\%$]
    \item Low Beta: $+3.9\%$ [95\% CI: $+1.2\%$ to $+6.5\%$]
    \item Theta: $+2.2\%$ [95\% CI: $-2.4\%$ to $+6.9\%$]
    \item Delta: $-12.1\%$ [95\% CI: $-22.0\%$ to $-2.5\%$]
\end{itemize}

The gamma confidence interval does not overlap with any other band, confirming that gamma's uniquely strong $\phisym^n$ alignment is not attributable to sampling variability.

\textbf{Position hierarchy analysis.} Beyond magnitude differences, we examined whether the predicted enrichment ordering (1° noble $>$ attractor $>$ 2° noble) is preserved within each band. \textit{Notation:} \checkmark\ = correct ordering; $\triangle$ = partial match; $\times$ = violated.

\begin{itemize}
    \item Gamma: $+144.8\% > +82.5\% > +6.1\%$ \checkmark\ (correct ordering)
    \item High Beta: $+8.8\% \approx +8.6\% \approx +8.3\%$ \checkmark\ (essentially tied but correct)
    \item Alpha: $+4.2\% > +2.8\% > +0.2\%$ \checkmark\ (correct ordering)
    \item Theta: $+2.2\% > -6.2\% > -16.7\%$ \checkmark\ (correct ordering)
    \item Low Beta: $+3.9\% > -4.6\%$, but $-4.6\% < -1.6\%$ $\triangle$ (partial)
    \item Delta: $-12.1\% < +58.6\%$ $\times$ (anomalous: attractor dominates over 1° noble)
\end{itemize}

Notably, 4 of 6 bands preserve the predicted hierarchy despite 10--50$\times$ magnitude differences. This supports underlying universality: the $\phisym^n$ architecture exists across bands with frequency-dependent \textit{expression strength}, not merely as a gamma phenomenon with incidental patterns elsewhere.

\begin{figure}[H]
    \centering
    \includegraphics[width=1\linewidth]{phi_band_heterogeneity.png}
    \caption{Inter-band heterogeneity analysis. \textbf{A.} Enrichment by band and position. Gamma (rightmost group) shows dramatically stronger effects at attractor ($+82.5\%$) and 1° noble ($+144.8\%$) positions compared to all other bands. \textbf{B.} Forest plot of 1° noble enrichment with 95\% bootstrap CIs. Gamma (green) is clearly separated from all other bands (blue), with non-overlapping confidence intervals. The gamma 1° noble enrichment is approximately 16$\times$ larger than the next highest band (high beta).}
    \label{fig:band_heterogeneity}
\end{figure}

\begin{criticalpoint}
\textbf{Interpretation: Universal hierarchy with frequency-dependent expression.} Gamma shows uniquely strong $\phisym^n$ alignment ($p < 10^{-50}$ vs. all other bands), but the predicted position hierarchy (1° noble $>$ attractor $>$ 2° noble) is preserved in 4 of 6 bands. This pattern---universal \textbf{ordering} with variable \textbf{magnitude}---suggests that different frequency bands implement the same $\phisym^n$ architecture with \textbf{frequency-dependent expression strength}.

Gamma's dominance is not anomalous but \textbf{mechanistically required}: temporal binding and feature integration demand precise phase relationships within $\sim$25 ms windows (matching GABA-A receptor kinetics), and the $\phisym^n$ architecture's anti-mode-locking property ensures oscillator independence essential for selective binding (see Section~\ref{sec:theoretical_interpretation} for mechanistic details). Lower-frequency bands serve functions---gating, memory consolidation, homeostatic regulation---that tolerate or benefit from frequency flexibility. The delta band's anomalous pattern (attractor dominance) may reflect fundamentally different organizational principles at slow timescales.

This is not merely a gamma phenomenon with incidental patterns elsewhere, but rather a shared organizational principle with \textbf{function-tuned expression strength}: operations requiring temporal precision show tight $\phisym^n$ adherence; operations integrating over longer timescales show looser constraints.
\end{criticalpoint}

\subsection{$\fzero$ Sensitivity Analysis}
\label{sec:f0_sensitivity}

To assess robustness to the choice of fundamental frequency, we computed the composite alignment metric across $\fzero$ values from 7.0 to 8.5 Hz in 0.05 Hz steps. The results reveal a characteristic plateau structure:

\begin{itemize}
    \item \textbf{Optimal range}: $\fzero \in [7.3, 7.9]$ Hz produces $>95\%$ of maximum alignment
    \item \textbf{Peak}: $\fzero = 7.60$ Hz yields maximum composite score
    \item \textbf{Canonical value}: $\fzero = 7.83$ Hz (literature Schumann value) produces 94\% of optimal
    \item \textbf{Degradation}: Below 7.0 Hz or above 8.2 Hz, alignment drops below 80\%
\end{itemize}

\begin{figure}[H]
    \centering
    \includegraphics[width=1\linewidth]{phi_f0_sensitivity.png}
    \caption{$\fzero$ sensitivity analysis. \textbf{Top:} Comprehensive alignment metric (ranking + magnitude) vs.\ base frequency, showing optimal at 7.60 Hz (green dashed) coinciding with geophysical Schumann measurements (black dotted). The canonical literature value (7.83 Hz, orange dotted) falls within the $>95\%$ plateau (shaded blue). \textbf{Bottom:} Position-type enrichment breakdown showing stable boundary depletion (green) and attractor enrichment (purple) across the plateau, with 1\degree\ noble enrichment (red) degrading sharply above 8.0 Hz as it collapses onto the attractor position.}
    \label{fig:f0_sensitivity}
\end{figure}

This 0.6 Hz plateau (7.3--7.9 Hz) is narrower than would be expected by chance---random grid positions show plateaus spanning 1.5--2.0 Hz. The plateau's centering on the measured Schumann fundamental is notable, though the statistical significance of this alignment requires careful examination.

\subsubsection{$\fzero$ Tolerance and Ratio Preservation}

The phase-shift permutation test ($p = 0.21$) reveals a $\sim$0.6 Hz tolerance plateau (7.3--7.9 Hz) within which alignment remains $>95\%$ of optimal (Figure~\ref{fig:f0_sensitivity}).

\begin{criticalpoint}
\textbf{This tolerance is expected, not a limitation.} The claim is not that the brain ``locks'' to a specific SR frequency---both SR ($\pm 0.3$ Hz diurnal) and neural $\fzero$ (individual differences, state fluctuations) are inherently variable. The $\phisym^n$ \textbf{ratio architecture} is the preserved quantity; absolute $\fzero$ can float within its natural range while the boundary-attractor structure remains intact.
\end{criticalpoint}

The 0.6 Hz tolerance plateau is consistent with ratio preservation in variable systems. The SIE data \citep{lacy_sie} show the same pattern at the event level: individual harmonic frequencies vary independently ($r \approx 0$ between SR1 and SR3), yet \textit{ratios} maintain $<1\%$ precision. At the population level, $\fzero$ can vary within $\pm 0.3$ Hz while preserving the boundary-attractor architecture.

The \textbf{open question} is not whether frequency-locking exists (it was never claimed), but \textit{why} $\fzero$ tends toward the SR range ($\sim$7.5 Hz):

\begin{enumerate}
    \item \textbf{Coincidental convergence:} $\phisym^n$ scaling arises from neural constraints (GABAergic time constants, thalamocortical loop delays) that independently center near 7.6 Hz.

    \item \textbf{Evolutionary/dynamic coupling:} Neural systems evolved or adapt to SR as a reference, maintaining $\phisym^n$ ratio organization across natural SR variability.
\end{enumerate}

Both hypotheses explain all observed data. The independent convergence of $\fzero$ from geophysical monitoring and SIE events (mean SR1 = 7.63 Hz; \citealt{lacy_sie}) is consistent with---but does not prove---coupling. Distinguishing these hypotheses requires experimental manipulation (e.g., SR-shielded vs.\ SR-exposed environments).

\subsection{Hierarchical Validation}

To address potential pseudoreplication from treating 244,955 peaks as independent observations, we performed hierarchical analyses aggregating to session-level and subject-level units of inference (see \S\ref{sec:hierarchical_methods}).

\subsubsection{Session-Level Statistics}

Computing one enrichment score per session (mean $n = 253$ peaks/session) and testing across $n = 968$ sessions yielded results consistent with peak-pooled estimates (Table~\ref{tab:hierarchical}).

\begin{table}[H]
\centering
\caption{Pooled vs. Hierarchical Analysis Comparison}
\label{tab:hierarchical}
\begin{tabular}{@{}lccc@{}}
\toprule
Metric & Peak-Pooled ($N$=244,955) & Session-Level ($N$=968) & Mixed Model \\
\midrule
\multicolumn{4}{l}{\textit{Attractor Enrichment}} \\
Mean & $+21.0\%$ & $+21.1\%$ & $+21.1\%$ \\
95\% CI & $[+20.5, +21.5]$ & $[+19.4, +22.9]$ & $[+19.4, +22.9]$ \\
Cohen's $d$ & --- & $0.80$ & --- \\
$p$-value & $< 0.0001$ & $< 0.0001$ & $< 0.0001$ \\
\midrule
\multicolumn{4}{l}{\textit{Boundary Enrichment}} \\
Mean & $-18.0\%$ & $-18.1\%$ & $-18.1\%$ \\
95\% CI & $[-18.4, -17.5]$ & $[-19.7, -16.4]$ & $[-19.7, -16.4]$ \\
Cohen's $d$ & --- & $-0.76$ & --- \\
$p$-value & $< 0.0001$ & $< 0.0001$ & $< 0.0001$ \\
\midrule
\multicolumn{4}{l}{\textit{Model Diagnostics}} \\
ICC (subjects) & --- & --- & 0.082 \\
Sessions/Subject & --- & 10.0 (mean) & --- \\
\bottomrule
\end{tabular}
\end{table}

\begin{figure}[H]
\centering
\includegraphics[width=0.85\textwidth]{phi_hierarchical_comparison.png}
\caption{Pooled vs. hierarchical analysis comparison. Bar plots show mean enrichment with 95\% confidence intervals for attractor (left) and boundary (right) positions across four estimation methods: peak-pooled estimates, session-level $t$-test, session-level bootstrap, and mixed-effects model with subject random intercept. All methods yield consistent estimates, confirming that effects are not inflated by pseudoreplication.}
\label{fig:hierarchical_comparison}
\end{figure}

Session-level $t$-tests confirmed both attractor enrichment ($t(967) = 24.5$, $p < 0.0001$, $d = 0.80$) and boundary depletion ($t(967) = -23.5$, $p < 0.0001$, $d = -0.76$). The paired test confirmed attractor $>$ boundary within sessions ($t(967) = 27.3$, $p < 0.0001$, $d = 0.89$), with 83.3\% of sessions showing positive paired differences (attractor minus boundary $> 0$).\footnote{This 83.3\% reflects the stricter paired-difference criterion; the 99.1\% reported in Section~4.5 reflects sessions where aggregate $E_{\text{attractor}} > E_{\text{boundary}}$. The difference arises because the paired test is more sensitive to within-session variance.}

\begin{figure}[H]
\centering
\includegraphics[width=\textwidth]{phi_session_level_enrichment_1.png}
\caption{Session-level enrichment distributions. Each panel shows the distribution of per-session enrichment scores ($n = 968$ sessions) for each lattice position type. Effect sizes (Cohen's $d$) and one-sample $t$-test $p$-values are shown. Boundary positions show consistent depletion ($d = -0.76$), attractor positions show consistent elevation ($d = 0.80$), and 1° Noble positions show the strongest effect ($d = 1.55$). The alignment score (attractor minus boundary) confirms the predicted ordering across sessions ($d = 1.00$).}
\label{fig:session_enrichment}
\end{figure}

\subsubsection{Mixed-Effects Modeling}

The linear mixed-effects model with subject random intercept yielded identical point estimates to session-level analyses (attractor: $\mu = +21.1\%$, 95\% CI: $[+19.4, +22.9]$, $p < 0.0001$). The low intraclass correlation (ICC $= 0.082$) indicates that only 8.2\% of variance in enrichment scores lies between subjects, with 91.8\% attributable to within-subject session variability. This suggests that $\phisym^n$ lattice alignment is consistent across individuals rather than driven by subject-specific characteristics.

\subsubsection{Session-Level Permutation Test}

Sign-flipping permutation at the session level (10,000 permutations) yielded $p < 0.0001$ for attractor enrichment, confirming that the effect is not an artifact of pseudoreplication.

\begin{criticalpoint}
\textbf{Hierarchical validation confirms key findings.} Session-level effect sizes (Cohen's $d = 0.80$ for attractor, $d = -0.76$ for boundary) remain large when each of 968 sessions contributes equally regardless of peak count. The low ICC (0.082) indicates minimal subject clustering, supporting the robustness of $\phisym^n$ alignment across individuals.
\end{criticalpoint}

\subsection{Cross-Device Validation}
\label{sec:cross_device}

To test whether the $\phisym^n$ architecture generalizes beyond Emotiv hardware, we analyzed 17 sessions recorded with the Interaxon MUSE headband (4 frontal/temporal channels) yielding 1,251 additional peaks. Despite different electrode placements, sensor technology, and fewer channels, MUSE data showed the same fundamental $\phisym^n$ organization:

\begin{itemize}
    \item \textbf{Position ordering preserved:} The theoretical hierarchy (1° Noble $>$ Attractor $>$ 2° Noble $>$ Boundary) was maintained
    \item \textbf{Boundary depletion:} $-15.8\%$ (vs.\ $-18.3\%$ Emotiv)
    \item \textbf{Attractor enrichment:} $+19.4\%$ (vs.\ $+22.2\%$ Emotiv)
    \item \textbf{1° Noble enrichment:} $+36.7\%$ (vs.\ $+39.3\%$ Emotiv)
\end{itemize}

\textbf{Band-specific differences.} While overall $\phisym^n$ structure was preserved, MUSE showed a distinct band profile compared to Emotiv:

\begin{itemize}
    \item \textbf{Gamma:} MUSE showed $+22.4\%$ enrichment at 1° Noble (vs.\ $+144.8\%$ Emotiv)---substantially weaker gamma effects
    \item \textbf{Alpha:} MUSE showed $+73.1\%$ enrichment (vs.\ $+4.2\%$ Emotiv)---much stronger alpha effects
    \item \textbf{Beta:} MUSE showed $+49.4\%$ enrichment (vs.\ $+5.9\%$ Emotiv)---stronger beta effects
\end{itemize}

These band differences likely reflect electrode placement: MUSE's frontal-temporal montage (AF7, AF8, TP9, TP10) captures different cortical sources than Emotiv's broader coverage. Frontal channels may emphasize alpha/beta rhythms while underrepresenting occipital-dominant gamma. Critically, the \textit{underlying} $\phisym^n$ architecture is preserved across devices---only the band-specific expression strength varies with electrode configuration.

\begin{keyinsight}
The $\phisym^n$ frequency architecture is \textit{device-independent}: both Emotiv and MUSE show the predicted boundary-attractor-noble hierarchy. Band-specific enrichment magnitudes vary with electrode placement, but the fundamental ratio structure is universal.
\end{keyinsight}

\subsection{Summary of Results}

\begin{table}[H]
\centering
\caption{Summary: Predictions vs. Observations}
\label{tab:summary}
\begin{tabular}{@{}p{5cm}p{4cm}p{4cm}c@{}}
\toprule
Prediction & Expected & Observed & Match \\
\midrule
\multicolumn{4}{l}{\textit{Position-type enrichment}} \\
Boundaries depleted & Low counts at integer $n$ & $-18\%$ enrichment & $\checkmark$ \\
Attractors elevated & High counts at half-integer $n$ & $+21\%$ enrichment & $\checkmark$ \\
1° Noble most enriched & Maximum at $n + 0.618$ & $+39\%$ enrichment & $\checkmark$ \\
2° Noble modestly enriched & Modest at $n + 0.382$ & $+2\%$ enrichment & $\checkmark$ \\
\midrule
\multicolumn{4}{l}{\textit{Frequency alignment}} \\
Alpha peak at $\phisym^{0.5}$ & 9.67 Hz & 9.8 $\pm$ 0.3 Hz & $\checkmark$ \\
$\beta/\gamma$ trough at $\phisym^3$ & 32.19 Hz & 32.4 $\pm$ 0.8 Hz & $\checkmark$ \\
Gamma peak at $\phisym^{3.5}$ & 40.95 Hz & 41.2 $\pm$ 0.6 Hz & $\checkmark$ \\
\midrule
\multicolumn{4}{l}{\textit{Statistical validation}} \\
Uniform frequency test & $\phisym^n$ structure in peaks & $p < 0.0001$ & $\checkmark$ \\
Lattice non-uniformity & Non-uniform within bands & $\chi^2 = 9{,}969$, $p < 10^{-100}$ & $\checkmark$ \\
$\phisym$ outperforms alternatives & Best alignment + correct ordering & $p < 0.001$ vs. all 5 & $\checkmark$ \\
$\fzero$ plateau includes SR value & 7.6 Hz in optimal range & 7.3--7.9 Hz ($>95\%$ optimal) & $\checkmark$ \\
\midrule
\multicolumn{4}{l}{\textit{Anchoring tolerance}} \\
Adaptive tolerance around SR & 0.6 Hz tolerance plateau & Matches SR variability range & $\checkmark$* \\
\midrule
\multicolumn{4}{l}{\textit{Robustness}} \\
Session-level consistency & Attractor $>$ boundary & 99.1\% (959/968 sessions) & $\checkmark$ \\
Effect size & Large difference & Cohen's $d \approx 0.8$--$1.0$ (session-level) & $\checkmark$ \\
Gamma strongest alignment & Highest $z$-scores in gamma & $z = +69.0$ (1° noble) & $\checkmark$ \\
\bottomrule
\end{tabular}
\begin{flushleft}
\small
*Consistent with adaptive coupling interpretation; distinguishing from coincidence requires experimental manipulation.
\end{flushleft}
\end{table}

%==============================================================================
\section{Discussion}
%==============================================================================

\subsection{Summary of Findings}

Analysis of 244,955 FOOOF-detected oscillatory peaks confirms that human EEG frequency organization follows golden ratio architecture:

\begin{enumerate}
    \item Peak frequencies show systematic depletion at integer $\phisym^n$ positions (boundaries: $-18\%$) and accumulation at half-integer positions (attractors: $+21\%$) and noble positions ($+39\%$)
    
    \item The optimal fundamental frequency ($\fzero = 7.6$ Hz) matches independently measured Schumann Resonance values
    
    \item The golden ratio outperforms all tested alternative scaling factors
    
    \item The pattern is consistent across sessions, subjects, contexts, and devices
    
    \item Predictions generated \textit{a priori} from geophysical data align with neural observations within 0.2 Hz
\end{enumerate}

\subsection{Evidence Strength Stratification}

Before detailed discussion, we stratify our claims by evidential support to help readers calibrate expectations:

\textbf{Strong evidence:}
\begin{itemize}
    \item Boundary depletion at $\phisym^3$ (32 Hz): $-18\%$, highly consistent
    \item Attractor enrichment at $\phisym^{3.5}$ (41 Hz): $+82.5\%$ in gamma
    \item $\phisym$ outperforms alternative scaling factors: $p < 0.001$ for all comparisons
    \item Uniform frequency permutation test: $p < 0.0001$
    \item Band heterogeneity: Gamma enrichment significantly exceeds all other bands ($p < 10^{-50}$ for all pairwise comparisons)
\end{itemize}

\textbf{Moderate evidence:}
\begin{itemize}
    \item Noble position enrichment pattern: Overall $+39\%$, variable across bands but consistently highest in gamma
    \item SR harmonic ratio precision: $\sim$0.6\% mean error, but mechanistically unexplained
    \item Frequency-dependent expression: Gamma shows tight $\phisym^n$ adherence; lower bands show weaker effects
\end{itemize}

\textbf{Preliminary/Unresolved:}
\begin{itemize}
    \item Why $\fzero$ near SR: Coincidental convergence from neural biophysics remains equally viable as coupling
    \item Lower-frequency $\phisym^n$ patterns: Weak effects may reflect IAF variability, edge effects, or genuine functional flexibility
\end{itemize}

\subsection{Addressing Potential Criticisms}

\subsubsection{``This is Just Curve Fitting''}

\textbf{Criticism:} Any framework with adjustable parameters can fit data post-hoc.

\textbf{Response:} The core equation has \textbf{zero fitted parameters} in the strict sense: $\fzero = 7.6$ Hz was established from geophysical monitoring before examining neural data, and $\phisym$ is a mathematical constant. Predictions were documented prior to analysis.

\textbf{However, we acknowledge analytical choices that affect results} (detailed in Section 2.4.3):
\begin{itemize}
    \item Using empirical $\fzero = 7.60$ Hz rather than theoretical $c/2\pi r = 7.49$ Hz (1.5\% difference)
    \item Window width of $\pm 0.5$ Hz for position enrichment
    \item Including noble positions (0.382, 0.618) alongside boundaries and attractors
\end{itemize}

These are not \textit{fitted} to maximize agreement---they were fixed before analysis based on independent rationale---but they do represent degrees of freedom that should be acknowledged. The key evidence against post-hoc fitting is: (a) $\fzero$ convergence from two independent sources (geophysical monitoring and SIE events: 7.63 Hz), and (b) results are qualitatively robust to reasonable variations in these choices.

\subsubsection{``The Pattern Could Be a FOOOF Artifact''}

\textbf{Criticism:} FOOOF parameters might artificially create peaks at certain frequencies.

\textbf{Response:} FOOOF is designed to detect oscillatory peaks above the aperiodic background without imposing frequency preferences. The algorithm has been validated extensively \citep{donoghue2020} and does not contain $\phisym$-related parameters.

We verified robustness by varying FOOOF parameters:
\begin{itemize}
    \item Peak width limits (0.5--12 Hz range): Pattern persists
    \item Peak threshold (1.5--3.0 SD): Pattern persists
    \item Maximum peaks (4--8): Pattern persists
\end{itemize}

The boundary-attractor pattern is not sensitive to reasonable parameter variation.

\subsubsection{``Consumer-Grade EEG is Unreliable''}

\textbf{Criticism:} Consumer devices may introduce artifacts not present in research-grade systems.

\textbf{Response:} Multiple lines of evidence address this concern:

\begin{enumerate}
    \item \textbf{Cross-device consistency:} The $\phisym^n$ pattern is consistent across three device families: Emotiv EPOC X (14-channel, saline sensors), Emotiv Insight (5-channel, semi-dry sensors), and Interaxon MUSE (4-channel, dry sensors). Each uses different sensor technologies, electrode placements, and amplifier designs, yet all show the predicted boundary-attractor-noble hierarchy (Section~\ref{sec:cross_device})

    \item The companion SIE paper demonstrates cross-device convergence with $<1\%$ ratio error in harmonic relationships

    \item The frequencies involved (5--50 Hz) are well within all three devices' validated ranges

    \item Hardware-specific artifacts would not produce $\phisym^n$ organization---there is no mechanism by which consumer EEG would preferentially create golden ratio structure

    \item While MUSE shows different band-specific magnitudes (stronger alpha, weaker gamma) consistent with its frontal-temporal electrode placement, the underlying $\phisym^n$ ratio architecture is preserved
\end{enumerate}

Nevertheless, replication with research-grade systems (64+ channels, 1000+ Hz sampling) is warranted and explicitly invited.

\subsubsection{``Why Hasn't This Been Discovered Before?''}

\textbf{Criticism:} With decades of EEG research, such a fundamental pattern should have been noticed.

\textbf{Response:} Several factors obscured the pattern:

\begin{enumerate}
    \item \textbf{Traditional analysis averages across bands:} Computing ``alpha power'' by averaging 8--13 Hz obscures within-band structure and precise peak locations
    
    \item \textbf{Pre-defined bands:} Analyzing power in fixed ranges assumes bands are correct, preventing discovery that boundaries might be elsewhere
    
    \item \textbf{1/f confound:} Raw power spectra mix oscillatory peaks with aperiodic background; FOOOF separation is relatively recent (2020)
    
    \item \textbf{Looking for the wrong things:} Research focused on band power correlations with behavior, not on the mathematical organization of frequencies themselves
    
    \item \textbf{Scale of data:} Testing $\phisym^n$ alignment requires many peaks across wide frequency range; 244,955 peaks provides statistical power unavailable in typical studies
\end{enumerate}

\subsubsection{``The Schumann Correspondence Could Be Coincidence''}

\textbf{Criticism:} Two numbers in similar ranges (neural $\fzero \approx 7.6$ Hz, Schumann $\approx 7.6$ Hz) might match by chance.

\textbf{Response:} We agree this is possible and explicitly frame the finding as \textbf{correspondence, not causation}. However:

\begin{enumerate}
    \item The match is precise: not ``roughly 8 Hz'' but specifically 7.60 Hz in both cases
    
    \item The value 7.6 Hz is \textit{optimal} for neural alignment---confirmed by independent convergence from SIE data (7.63 Hz) and spectral analysis
    
    \item The Schumann value was measured independently by multiple observatories over multiple years
    
    \item If coincidental, it is a remarkable coincidence that the Earth's electromagnetic cavity frequency and the human brain's organizing frequency match to within 0.5\%
\end{enumerate}

Distinguishing coincidence from coupling requires experiments beyond this observational study (see Future Directions).

\subsubsection{``Golden Ratio Claims are Pseudoscience''}

\textbf{Criticism:} The golden ratio has been overclaimed in many contexts; this might be another case.

\textbf{Response:} We acknowledge historical overclaiming and have taken specific measures to avoid it:

\begin{enumerate}
    \item \textbf{Quantitative predictions:} We specify exact frequencies, not vague ``approximately phi'' claims
    
    \item \textbf{Statistical testing:} We compare $\phisym$ to five alternatives with proper statistics; $\phisym$ significantly outperforms all ($p < 0.001$) and is the only factor exhibiting the correct theoretical ordering
    
    \item \textbf{Mechanistic justification:} We explain \textit{why} $\phisym$ should be special (anti-commensurability, Fibonacci property), not just that it is
    
    \item \textbf{Falsifiable predictions:} The framework makes specific predictions (e.g., 1° noble at 10.23 Hz, not 10.5 Hz) that can be tested
    
    \item \textbf{No mysticism:} We invoke coupled oscillator dynamics, not numerology
\end{enumerate}

The question is empirical: does $\phisym$-based prediction match observation better than alternatives? The data answer yes.

\subsubsection{``Effects Are Inflated by Pseudoreplication''}

\textbf{Criticism:} Treating 244,955 peaks as independent observations inflates sample size and overstates effect sizes, since peaks are nested within channels within sessions within subjects.

\textbf{Response:} We addressed this concern through hierarchical analysis (Section~4.5):

\begin{enumerate}
    \item \textbf{Session-level inference:} Computing one enrichment score per session and testing across $n = 968$ sessions yielded effect sizes comparable to peak-pooled estimates (Cohen's $d = 0.80$ for attractor enrichment, $d = -0.76$ for boundary depletion). Confidence intervals widened appropriately (from $\pm 0.5\%$ to $\pm 1.7\%$) but effects remained highly significant ($p < 0.0001$).

    \item \textbf{Mixed-effects modeling:} Accounting for subject clustering via random intercepts confirmed the effect remained significant ($p < 0.0001$). The low intraclass correlation (ICC $= 0.082$) indicates that 91.8\% of variance is within-subject, suggesting the pattern is consistent across individuals rather than driven by a subset of subjects.

    \item \textbf{Session-level permutation:} Sign-flipping permutation testing at the session level (10,000 permutations) yielded $p < 0.0001$, confirming that effects are not artifacts of pseudoreplication.

    \item \textbf{Cross-session consistency:} 83.3\% of individual sessions show attractor $>$ boundary enrichment, demonstrating the pattern is not driven by a few outlier sessions.
\end{enumerate}

Key effects persist when each session contributes equally regardless of peak count. Session-level effect sizes remain ``large'' by conventional standards ($d > 0.8$), confirming that findings are robust to proper hierarchical accounting.

\subsection{Theoretical Interpretation}
\label{sec:theoretical_interpretation}

Having addressed methodological concerns, we now examine the theoretical basis for the observed $\phisym^n$ organization.

\subsubsection{Why Boundaries Are Depleted}

At integer $\phisym^n$ positions, the Fibonacci property creates a unique resonance condition:

\begin{equation}
f(n) = f(n-1) + f(n-2)
\end{equation}

This enables three-wave mixing: energy at $f(n)$ can transfer to $f(n-1)$ and $f(n-2)$ (or vice versa) through nonlinear interactions. Oscillations at boundary frequencies are therefore \textbf{unstable}---they act as transition points where energy redistributes across bands rather than accumulating.

\subsubsection{Why Attractors Are Elevated}

At half-integer positions ($n = k + 0.5$), oscillations are maximally distant (in log-frequency space) from boundaries on either side. The golden ratio's anti-commensurability means these frequencies cannot easily couple to other positions. Oscillations at attractors are \textbf{stable}---they resist perturbation and accumulate over time.

\subsubsection{Why Noble Positions Show Maximum Enrichment}

The 1° noble positions ($n = k + 0.618 = k + 1/\phisym$) inherit the golden ratio's anti-locking property at a finer scale. They represent the ``most golden'' positions within each band---maximally distant from simple rational relationships with both boundaries. This explains why the canonical alpha peak is near 10.2 Hz ($\phisym^{0.618}$) rather than 9.7 Hz ($\phisym^{0.5}$).

\subsubsection{The Segregation-Integration Balance}

The $\phisym^n$ architecture achieves a balance central to neural computation:

\begin{itemize}
    \item \textbf{Segregation:} Oscillations at different frequencies remain independent (resist mode-locking) because $\phisym$ is maximally irrational
    
    \item \textbf{Integration:} Energy can transfer between bands at boundaries because the Fibonacci property enables resonant coupling
\end{itemize}

No other scaling factor provides both properties. Integer ratios (2:1, 3:2) enable coupling but produce mode-locking. Other irrationals ($\sqrt{2}$, $e$) resist mode-locking but lack Fibonacci coupling. The golden ratio is unique in optimizing both.

\subsubsection{Convergence of Independent $\fzero$ Estimates}

The fundamental frequency $\fzero = 7.6$ Hz was established from two independent sources:

\begin{enumerate}
    \item \textbf{Geophysical monitoring}: Multi-year Schumann Resonance measurements at Tomsk Observatory yield $7.6 \pm 0.2$ Hz for the fundamental mode.

    \item \textbf{Transient neural events}: Mean SR1 frequency from 1,121 Schumann Ignition Events with valid harmonic measurements (of 1,366 total SIEs detected across 91 subjects) is $7.63 \pm 0.33$ Hz---consistent with the empirical $\fzero = 7.60$ Hz used throughout this analysis (see Section~2.3.4 for the relationship between theoretical and empirical values).
\end{enumerate}

This convergence from geophysical and transient neural data supports the framework's validity. The choice of $\fzero$ was not tuned to maximize alignment metrics but rather derived from independent geophysical measurements.

\begin{keyinsight}
The $\fzero = 7.6$ Hz value emerges independently from two measurement domains---atmospheric electromagnetics and transient high-coherence neural states---providing cross-modal validation of the Schumann-neural coupling hypothesis.
\end{keyinsight}

The $\fzero$ sensitivity analysis (Figure~\ref{fig:f0_sensitivity}) reveals a 0.6 Hz plateau (7.3--7.9 Hz) where alignment remains $>95\%$ of optimal. This plateau width explains the phase-shift permutation result ($p = 0.21$): approximately 20\% of random phase shifts land within the acceptable plateau. The $\phisym^n$ architecture appears tolerant to $\sim\pm0.3$ Hz variation in $\fzero$, consistent with biological variability in Schumann Resonance measurements ($\pm 0.2$ Hz diurnal variation) and individual differences in neural reference frequencies.

\subsubsection{Theoretical Advantage of the Golden Ratio}

The golden ratio's unique suitability for neural frequency organization derives from its mathematical properties rather than simple alignment metrics:

\begin{enumerate}
    \item \textbf{Maximal anti-commensurability}: $\phisym$ is the ``most irrational'' number (slowest-converging continued fraction), making frequencies at $\phisym^n$ ratios maximally resistant to mode-locking.

    \item \textbf{Fibonacci additivity}: Only $\phisym$ satisfies $\phisym^n = \phisym^{n-1} + \phisym^{n-2}$, enabling three-wave resonant coupling at boundary positions while maintaining independence at attractors.
\end{enumerate}

These properties predict that oscillations at attractor positions ($n + 0.5$) should be stable, while boundaries (integer $n$) serve as energy transfer points. The observed enrichment pattern (attractors: $+21\%$; boundaries: $-18\%$) is consistent with this theoretical framework.

\begin{keyinsight}
The theoretical argument for $\phisym$ is not that it maximizes some generic alignment metric, but that its unique mathematical properties---anti-commensurability and Fibonacci additivity---make it optimally suited for balancing neural segregation and integration.
\end{keyinsight}

\subsection{Relationship to Companion SIE Paper}

The present paper and its companion \citep{lacy_sie} investigate the same $\phisym^n$ phenomenon at complementary scales of analysis. Together, they provide convergent evidence from independent methodologies:

\begin{table}[H]
\centering
\caption{Complementary Evidence from Two Independent Analyses}
\label{tab:paper_relationship}
\begin{tabular}{@{}ll@{}}
\toprule
Paper 1: Schumann Ignition Events & Paper 2: Golden Ratio Architecture \\
\midrule
Transient high-coherence states & Continuous baseline organization \\
1,366 events (91 subjects) & 244,955 peaks (968 sessions) \\
$\phisym^n$ ratios with $<1\%$ error & $\phisym^n$ structure $\pm 0.2$ Hz \\
SR3/SR1 = 2.6223 ($+0.16\%$ from $\phisym^2$) & Boundary depletion $-18\%$, attractor enrichment $+21\%$ \\
State-independent ($p = 0.69$) & Session-consistent (99.1\% show pattern) \\
Discovery of phenomenon & Framework explaining phenomenon \\
Events are ``ignitions'' & Architecture is ``substrate'' \\
\bottomrule
\end{tabular}
\end{table}

\subsubsection{Mutual Reinforcement of Findings}

The two analyses reinforce each other in several important ways:

\textbf{1. Temporal scale complementarity.} The SIE paper identifies $\phisym^n$ organization in \textit{transient} events (5--15 seconds), while the present paper demonstrates the same organization in \textit{aggregate} peak distributions across entire sessions. If $\phisym^n$ were an artifact of either analysis method, we would not expect it to appear at both timescales. The convergence suggests a genuine organizational principle manifesting at multiple temporal scales---transiently during high-coherence episodes, and continuously in baseline spectral structure.

\textbf{2. Precision gradient.} The SIE events show extraordinarily tight ratio precision ($<1\%$ deviation from $\phisym^n$ predictions), while the aggregate architecture shows broader but still significant structure ($\pm 0.2$ Hz tolerance). This gradient is \textit{expected} if SIE events represent moments when the underlying $\phisym^n$ architecture is maximally expressed: signal averages over many events should reveal the attractor structure with high precision, while population-level aggregation introduces individual variability that broadens the distribution.

\textbf{3. Independent $\fzero$ convergence.} Both papers arrive at the same fundamental frequency through entirely different analytical paths:
\begin{itemize}
    \item \textbf{SIE analysis}: Mean SR1 frequency across 1,121 events with valid harmonic measurements = $7.63 \pm 0.33$ Hz
    \item \textbf{Spectral architecture}: Optimal alignment from geophysical monitoring and phase-shift analysis = $7.60 \pm 0.2$ Hz
\end{itemize}
This convergence from transient neural events and aggregate spectral structure strongly supports the validity of both analyses. If either value were methodologically biased, we would not expect agreement to within 0.03 Hz.

\subsubsection{The Independence-Convergence Paradox: Theoretical Integration}

The companion SIE paper reports a striking finding that demands theoretical explanation: individual harmonic frequencies vary independently across events ($r \approx 0$ between SR1 and SR3), yet their \textit{ratios} maintain $\phisym^n$ precision ($<1\%$ error). We term this the ``independence-convergence paradox.''

This paradox is not merely an interesting observation---it has profound implications for understanding $\phisym^n$ neural organization:

\textbf{What it rules out:}
\begin{itemize}
    \item \textbf{Simple proportional scaling}: If all harmonics scaled together (e.g., from a common master oscillator), SR1 and SR3 would be strongly correlated. The near-zero correlation ($r \approx 0$) rules out this mechanism.
    \item \textbf{Absolute frequency locking}: If the brain ``locked'' to specific Schumann frequencies, we would expect tight clustering around fixed values. Instead, individual frequencies float freely.
    \item \textbf{Statistical artifact}: Random frequency variation would not produce ratio precision---it would produce ratio noise matching frequency noise.
\end{itemize}

\textbf{What it implies:}
\begin{itemize}
    \item \textbf{Ratio-level constraint}: The neurally constrained quantity is not \textit{frequency} but \textit{frequency ratio}. Individual oscillators can drift in absolute frequency while maintaining precise ratio relationships.
    \item \textbf{Distributed implementation}: The ratio constraint must be implemented across the network, not at individual oscillator sites. This is consistent with the network-level topological mechanisms proposed in Section~5.7.2.
    \item \textbf{Attractor dynamics}: The paradox suggests that the dynamically stabilized variable is the ratio $f_i/f_j$, with individual frequencies as dependent variables that can fluctuate within the ratio constraint.
\end{itemize}

\textbf{Integration with the present framework:}

The boundary-attractor-noble hierarchy identified in this paper provides the \textit{structural substrate} within which the independence-convergence paradox operates. Consider the following synthesis:

\begin{enumerate}
    \item The $\phisym^n$ lattice defines a \textit{frequency landscape} with stable positions (attractors, nobles) and unstable positions (boundaries).

    \item During baseline activity, oscillatory peaks preferentially occupy stable positions, producing the aggregate enrichment/depletion pattern observed here (Table~\ref{tab:position_enrichment}).

    \item During SIE events, multiple oscillators transiently achieve high coherence. The ratio constraint ensures that even as individual frequencies vary (perhaps due to noise, state fluctuations, or individual differences), the \textit{relationship} between harmonics remains fixed at $\phisym^n$.

    \item The 0.6 Hz tolerance plateau (Section~\ref{sec:phase_shift}) is the population-level manifestation of this same principle: $\fzero$ can vary within a range while preserving ratio architecture.
\end{enumerate}

This synthesis explains why both papers observe $\phisym^n$ organization despite examining different phenomena: they are measuring the same underlying constraint at different levels of analysis. The SIE paper measures ratio precision in individual events; this paper measures the statistical signature of ratio constraints across the population.

\subsubsection{Substrate-Ignition Model}

We propose a ``substrate-ignition'' model integrating both sets of findings:

\begin{keyinsight}
The $\phisym^n$ architecture exists continuously as the baseline organization of EEG frequencies---a \textit{substrate} that constrains where oscillatory activity can stably occur. Schumann Ignition Events represent transient episodes when coherence amplifies across this substrate, producing the ultra-precise ratio relationships observed in the SIE analysis. The substrate enables the ignitions; the ignitions reveal the substrate's precision.
\end{keyinsight}

This model makes a testable prediction: during SIE events, the distribution of \textit{individual} harmonic frequencies (SR1, SR3, etc.) should cluster at the attractor and noble positions defined by the $\phisym^n$ lattice, not at arbitrary frequencies. Preliminary analysis supports this prediction (SIE SR1 mean = 7.63 Hz, corresponding to $\phisym^0$ boundary position), but systematic validation requires correlating SIE harmonic frequencies with the lattice coordinate framework developed here.

The state-independence of SIE properties (ANOVA: SR1 $p = 0.21$, SR3 $p = 0.93$, SR5 $p = 0.15$ across cognitive contexts) is consistent with this model: if $\phisym^n$ represents an intrinsic architectural constraint rather than a state-dependent phenomenon, SIE events should show similar ratio precision regardless of the cognitive state in which they occur. This is precisely what the companion paper observes.

\subsection{Principled Band Definitions}

Beyond explaining existing observations, the framework yields practical outputs: mathematically derived band definitions that resolve historical inconsistencies:

\begin{table}[H]
\centering
\caption{$\phisym^n$-Based EEG Band Definitions}
\label{tab:band_defs}
\begin{tabular}{@{}lcccc@{}}
\toprule
Band & Lower Boundary & Attractor & Upper Boundary & Width \\
\midrule
Delta & --- & --- & $\phisym^{-1}$ = 4.70 Hz & --- \\
Theta & 4.70 Hz & $\phisym^{-0.5}$ = 5.97 Hz & $\phisym^{0}$ = 7.60 Hz & 2.90 Hz \\
\textbf{Alpha} & 7.60 Hz & $\phisym^{0.5}$ = 9.67 Hz & $\phisym^{1}$ = 12.30 Hz & 4.70 Hz \\
Low Beta & 12.30 Hz & $\phisym^{1.5}$ = 15.64 Hz & $\phisym^{2}$ = 19.90 Hz & 7.60 Hz \\
High Beta & 19.90 Hz & $\phisym^{2.5}$ = 25.31 Hz & $\phisym^{3}$ = 32.19 Hz & 12.29 Hz \\
\textbf{Gamma} & 32.19 Hz & $\phisym^{3.5}$ = 40.95 Hz & $\phisym^{4}$ = 52.09 Hz & 19.90 Hz \\
\bottomrule
\end{tabular}
\end{table}

These definitions resolve historical inconsistencies by providing principled boundaries. Note that band widths follow Fibonacci progression: $2.90, 4.70, 7.60, 12.29, 19.90$ Hz (approximately 3, 5, 8, 13, 21 when rounded).

\subsection{Explanation of Canonical Frequencies}

The framework explains longstanding questions:

\textbf{Why does alpha peak at $\sim$10 Hz?}

The $\phisym^{0.5}$ attractor is at 9.67 Hz; the 1° noble is at 10.23 Hz. The alpha peak represents accumulation at these stable positions.

\textbf{Why is 40 Hz special for gamma?}

The $\phisym^{3.5}$ attractor is at 40.95 Hz. This is not coincidence---it reflects the intersection of two independent constraints:

\begin{enumerate}
    \item \textbf{Biophysical constraint}: GABA-A receptor-mediated inhibition has a decay constant of $\sim$25 ms, creating a natural refractory period that favors oscillations near 40 Hz \citep{buzsaki2006}. This timescale is optimal for the rapid inhibitory cycling required to generate precise temporal windows for synaptic integration.

    \item \textbf{Architectural constraint}: The $\phisym^n$ framework places the gamma attractor at 40.95 Hz, where the anti-mode-locking property of the golden ratio ensures that distributed gamma oscillators maintain independence rather than collapsing into trivial synchrony.
\end{enumerate}

The convergence of biophysical and architectural constraints at $\sim$40 Hz explains why this frequency dominates binding, attention, and consciousness research \citep{gray1989, iaccarino2016}. Feature integration requires oscillators that are (a) fast enough to create narrow temporal windows, and (b) resistant to spurious synchronization. Only the $\phisym^{3.5}$ position satisfies both requirements.

\textbf{Why is beta ``broad'' and amorphous?}

Beta spans two $\phisym$-octaves ($\phisym^1$ to $\phisym^3$), containing two sub-bands (low and high beta) with distinct attractors. Its apparent breadth reflects this dual structure.

\textbf{Why is the beta-gamma boundary controversial (25--35 Hz)?}

The true boundary is at $\phisym^3 = 32.19$ Hz. Published values ranging from 25 to 35 Hz represent variation around this position.

\subsection{The Schumann Correspondence: Three Hypotheses}

The $\phisym^n$ ratio architecture is established; the open question is \textit{why} $\fzero$ tends toward the SR range ($\sim$7.5 Hz). Three hypotheses deserve consideration:

\subsubsection{Hypothesis 1: Evolutionary Tuning}

Nervous systems evolved over hundreds of millions of years in an environment permeated by Schumann Resonance fields. Neural oscillation frequencies may have been tuned by natural selection to match environmental electromagnetic rhythms.

\textbf{Testable prediction:} Species with longer evolutionary exposure to SR should show better $\phisym^n$ alignment than recently evolved species.

\subsubsection{Hypothesis 2: Biophysical Convergence via Evolutionary Optimization}

The $\sim$7.5 Hz fundamental may represent a convergent solution to independent constraints:

\begin{itemize}
    \item GABA-A receptor decay ($\tau \sim 25$ ms $\rightarrow f \sim 40$ Hz)
    \item GABA-B receptor decay ($\tau \sim 150$ ms $\rightarrow f \sim 7$ Hz)
    \item Thalamocortical loop delay ($\sim$100 ms $\rightarrow f \sim 10$ Hz)
    \item Membrane time constant ($\tau_m \sim 20$ ms $\rightarrow f \sim 50$ Hz)
\end{itemize}

Multiple timescales interact to favor $\sim$7.5 Hz as a fundamental.

Critically, these biophysical timescales are \textit{themselves} evolutionary adaptations, optimized over hundreds of millions of years for efficient neural computation. GABA receptor kinetics, thalamocortical conduction delays, and membrane time constants have been independently tuned by natural selection to support cognitive operations. The $\phisym^n$ architecture emerges as a \textit{secondary consequence}: when evolution independently optimizes multiple oscillatory systems, their interaction produces frequency ratios that naturally scale according to $\phisym^n$.

This represents a \textbf{stronger claim than direct SR coupling}: the brain does not ``track'' Schumann Resonance---it independently converged on similar values because $\phisym^n$ scaling is the mathematically optimal solution for coupled oscillator systems, and biophysical constraints center $\fzero$ near 7--8 Hz. The 0.6 Hz tolerance plateau (Section 4.5) reflects this intrinsic optimization rather than adaptive tracking of external fields.

The $\phisym^n$ architecture may thus be a \textbf{universal solution} that any sufficiently complex neural system would discover. If evolution optimizes individual biophysical timescales independently, the resulting oscillation frequencies will exhibit $\phisym^n$ ratio relationships as an emergent property---not because the brain ``knows'' about the golden ratio, but because $\phisym^n$ scaling uniquely satisfies the competing demands of segregation (anti-mode-locking) and integration (Fibonacci coupling).

\begin{keyinsight}
Hypotheses 1 and 2 are not mutually exclusive---both are evolutionary processes. The difference is whether SR is the \textit{target} of adaptation (Hypothesis 1) or a \textit{coincidental match} arising from independent optimization within biophysical constraints (Hypothesis 2). Hypothesis 2 is arguably more parsimonious: it requires no external coupling mechanism, only optimization within physical constraints.
\end{keyinsight}

\textbf{Testable prediction:} Artificially varying biophysical parameters (e.g., GABAergic modulation) should shift $\fzero$.

\subsubsection{Hypothesis 3: Direct Coupling}

Neural tissue may respond directly to SR fields through:

\begin{itemize}
    \item Stochastic resonance amplification of weak fields
    \item Coherent quantum effects
    \item Extremely sensitive detection mechanisms
\end{itemize}

\textbf{Testable prediction:} SR-shielded environments should show degraded $\phisym^n$ alignment; SR-enhanced environments should show improved alignment.

\begin{criticalpoint}
The current study demonstrates correspondence, not causation. Distinguishing these hypotheses requires experimental manipulation beyond observational analysis. We make no claim about mechanism.
\end{criticalpoint}

\subsection{Limitations}

\subsubsection{Device Limitations}

Data were collected with consumer-grade EEG (Emotiv EPOC/Insight, 5--14 channels; Interaxon MUSE, 4 channels; all at 128 Hz). While validated for research \citep{badcock2013, krigolson2017} and showing cross-device consistency across three hardware platforms in our data, replication with high-density research-grade systems (64+ channels, 1000+ Hz) is warranted.

\subsubsection{Population Limitations}

Participants were primarily healthy adults. The developmental trajectory of $\phisym^n$ architecture (infancy through aging) is unknown. Clinical populations may show altered organization.

\subsubsection{FOOOF Limitations}

FOOOF assumes oscillatory peaks are Gaussian and aperiodic background follows a power law. Deviations from these assumptions could affect peak detection. Alternative spectral methods (wavelets, IRASA) should be tested.

\subsubsection{Frequency Range Limitations}

Our analysis covered 1--48 Hz. Delta ($<4$ Hz) was undersampled due to high-pass filtering. High gamma ($>50$ Hz) was excluded due to Nyquist limitations. Extension to broader ranges requires different recording parameters.

\subsubsection{Causal Limitations}

The SR-neural correspondence is observational. We cannot distinguish evolutionary tuning, biophysical convergence, or direct coupling without experimental manipulation.

\subsubsection{Ratio Preservation Across $\fzero$ Variability}

\textbf{The $\phisym^n$ ratio architecture is robustly established.} The uniform frequency test ($p < 0.0001$) demonstrates significant $\phisym^n$ structure in EEG peak frequencies---the core finding of this paper. The phase-shift test ($p = 0.21$) reveals a $\sim$0.6 Hz tolerance plateau, which is \textbf{expected} given that both SR and neural $\fzero$ are naturally variable systems. Ratios, not absolute frequencies, are the preserved quantity.

The phase-shift result is consistent with ratio preservation in two variable systems. Both Schumann Resonance ($\pm 0.3$ Hz diurnal variation) and neural $\fzero$ (individual differences, state fluctuations) vary naturally. The 0.6 Hz plateau encompasses this variability. Frequency-locking to a precise value was never claimed; the claim is that \textit{ratio relationships} (boundary depletion, attractor enrichment) persist across the natural range of $\fzero$ variation.

This prediction is confirmed by the SIE data \citep{lacy_sie}: individual harmonic frequencies vary independently ($r \approx 0$ between SR1 and SR3), yet \textit{ratios} maintain $<1\%$ precision. The population-level $\fzero$ tolerance mirrors this pattern at the architectural level.

The \textbf{open question} is why $\fzero$ tends toward the SR range ($\sim$7.5 Hz): (1) coincidental convergence from neural biophysics, or (2) evolutionary/dynamic coupling to SR. Both hypotheses explain all observed data; distinguishing them requires experimental manipulation (e.g., SR-shielded environments).

\subsubsection{Mechanistic Gap: How Are Ratio Constraints Implemented?}

Section~5.5.2 provides theoretical integration of the independence-convergence paradox from the companion SIE paper, explaining \textit{what} is constrained ($\phisym^n$ ratios rather than absolute frequencies) and how this integrates with the substrate-ignition model. However, a critical mechanistic question remains: \textit{how} does the neural system implement ratio-level constraints while allowing individual frequencies to vary independently?

We propose three candidate mechanisms requiring future investigation:

\begin{enumerate}
    \item \textbf{Homeostatic frequency regulation}: Independent feedback loops maintain each harmonic at its $\phisym^n$-constrained position. The SIE data reveal a compensatory error mechanism ($r = -0.314$ between SR3/SR1 error and SR5/SR3 error): when one ratio deviates above $\phisym^n$, the adjacent ratio compensatorily deviates below, preserving overall precision. This suggests coordinated but not synchronized regulation across frequency bands.

    \item \textbf{Network-level topological constraints}: The $\phisym^n$ architecture may emerge from cortical network geometry (e.g., connectivity distributions, path-length statistics, conduction delays) that constrains frequency ratios without requiring explicit phase-locking. Cross-device consistency supports this: three device families with different electrode placements all show the predicted pattern (Section~\ref{sec:cross_device}), suggesting the constraint operates at network rather than sensor level.

    \item \textbf{Attractor dynamics in ratio space}: Rather than individual frequencies being controlled, the \textit{ratio} $f_i/f_j$ may be the dynamically stabilized quantity, with individual frequencies as free variables that fluctuate within the ratio constraint. This is analogous to a pendulum maintaining constant period despite amplitude variations---period is the attractor, not amplitude.
\end{enumerate}

Distinguishing these mechanisms requires time-resolved analysis of how harmonic frequencies covary within individual SIE events and whether ratio deviations trigger compensatory adjustments. The independence-convergence paradox represents a key mechanistic gap: both papers demonstrate \textit{what} is constrained; the neural implementation remains an open question.

\subsubsection{Band-Specific Variation}

As detailed in Section~4.4.7, the $\phisym^n$ alignment pattern varies significantly across frequency bands: gamma shows uniquely strong 1° noble enrichment ($+144.8\%$, $p < 10^{-50}$ vs.\ all other bands), while lower-frequency bands show weaker effects. However, magnitude heterogeneity does not imply absence of underlying structure---the predicted ordering (1° noble $>$ attractor $>$ 2° noble) is preserved in 4 of 6 bands despite 10--50$\times$ magnitude differences.

The delta band represents an exception, showing inverted hierarchy (attractor dominance over 1° noble). This may reflect edge effects near infraslow oscillations, distinct functional organization for slow cortical rhythms, or fundamentally different organizational principles at these timescales.

This heterogeneity admits two interpretations:

\begin{enumerate}
    \item \textbf{Methodological interpretation}: Higher-frequency oscillations have narrower spectral precision relative to bandwidth, producing sharper lattice coordinate distributions. Lower frequencies suffer more from individual frequency variability (IAF $\pm$1 Hz smears alpha) and edge effects (theta contains $\fzero$ at its boundary).

    \item \textbf{Functional interpretation}: Different bands implement $\phisym^n$ constraints for different purposes, reflecting their computational roles:

    \begin{itemize}
        \item \textbf{Gamma ($\phisym^3$--$\phisym^4$)}: Subserves temporal binding, feature integration, and conscious perception. These functions require precise phase relationships within $\sim$25 ms windows (matching GABA-A kinetics). Strict $\phisym^n$ adherence ensures that distributed gamma oscillators remain independent rather than collapsing into trivial synchrony---a requirement for selective feature binding.

        \item \textbf{Alpha ($\phisym^0$--$\phisym^1$)}: Functions primarily as a gating mechanism, suppressing irrelevant inputs and modulating cortical excitability. Gating operations are inherently flexible: they must accommodate variable input patterns. Looser $\phisym^n$ constraints may actually benefit alpha's functional role.

        \item \textbf{Theta ($\phisym^{-1}$--$\phisym^0$)}: Supports memory encoding and hippocampal-cortical communication. These slower operations integrate information over longer timescales and may tolerate broader frequency variability.

        \item \textbf{Delta ($<\phisym^{-1}$)}: Serves homeostatic and slow-wave sleep functions operating on timescales of seconds. The anomalous attractor-dominance pattern may reflect fundamentally different organizational principles at these slow timescales.
    \end{itemize}

    This functional gradient---from strict gamma compliance to flexible delta organization---suggests that $\phisym^n$ constraints are \textbf{tuned to computational requirements}: operations requiring precise temporal coordination show tight adherence; operations integrating over longer timescales show looser constraints.
\end{enumerate}

We favor a combination: band-specific \textit{functional requirements} interact with \textit{methodological factors} to produce the observed heterogeneity. Our claim is that the $\phisym^n$ architecture manifests as a \textbf{universal position hierarchy with frequency-dependent expression strength}---gamma shows uniquely strong expression, but the organizational principle (1° noble $>$ attractor $>$ 2° noble) is preserved across most bands. The delta band represents an exception requiring further investigation.

\subsection{Future Directions}

\subsubsection{Replication with Research-Grade EEG}

Independent replication using 64-channel systems (e.g., BioSemi, EGI) with 500+ Hz sampling is the most immediate priority. The prediction is specific: boundaries at integer $\phisym^n \times 7.6$ Hz, attractors at half-integers, nobles at $k + 0.618$.

\subsubsection{Cross-Species Comparison}

If $\phisym^n$ organization reflects evolutionary tuning to SR, species with similar neural architectures (mammals) should show similar patterns, while species with different architectures (invertebrates) might differ. Preliminary octopus EEG data show intriguing patterns warranting systematic investigation.

\subsubsection{Developmental Trajectory}

Recording from infants through elderly could reveal whether $\phisym^n$ architecture is present at birth, emerges developmentally, or degrades with aging.

\subsubsection{Clinical Applications}

Deviation from $\phisym^n$ organization may serve as a biomarker. Conditions involving oscillatory dysfunction (epilepsy, schizophrenia, Alzheimer's) could show altered boundary-attractor patterns.

\subsubsection{Concurrent EEG-Magnetometer Recording}

Recording EEG simultaneously with SR monitoring could test for direct coupling. If neural oscillations track SR fluctuations in real-time, direct coupling is supported; if they don't, evolutionary/biophysical explanations are favored.

\subsubsection{Stimulation at $\phisym^n$ Frequencies}

Transcranial alternating current stimulation (tACS) or sensory entrainment at $\phisym^n$ attractor frequencies might produce stronger effects than at arbitrary frequencies. This could have therapeutic applications.

%==============================================================================
\section{Conclusions}
%==============================================================================

Analysis of 244,955 oscillatory peaks across 968 sessions recorded with three consumer-grade EEG platforms (Emotiv EPOC X, Emotiv Insight, Interaxon MUSE) reveals that human EEG frequency organization follows golden ratio architecture:

\begin{keyequation}
\begin{align}
f(n) &= \fzero \times \phisym^n = \frac{c}{2\pi r} \times \phisym^n \quad \text{(Hz)} \\[0.3cm]
\omega(n) &= \frac{c}{r} \times \phisym^n \quad \text{(rad/s)}
\end{align}

\vspace{0.2cm}
\noindent where $c$ is the speed of light, $r$ is Earth's radius, and $\phisym$ is the golden ratio. The fundamental frequency $\fzero = c/(2\pi r) \approx 7.5$ Hz emerges directly from planetary geometry.
\end{keyequation}

\noindent The $\phisym^n$ \textit{ratio architecture} is robustly established ($p < 0.0001$). The phase-shift test reveals a $\sim$0.6 Hz tolerance plateau around $\fzero$, which is expected given that both SR and neural oscillations are variable systems---\textbf{ratios, not absolute frequencies, are the deep invariant}. The open question is \textit{why} $\fzero$ tends toward the SR range: coincidental convergence from neural biophysics, or evolutionary/dynamic coupling.

\noindent Integer $n$ positions function as boundaries (depleted, $-18\%$); half-integer positions function as attractors (enriched, $+21\%$); primary noble positions show maximum enrichment ($+39\%$).

This framework:

\begin{enumerate}
    \item Provides the first mathematically principled basis for EEG band definitions
    \item Explains why canonical rhythms occur at their characteristic frequencies
    \item Contains zero \textit{fitted} parameters in the core equation, though analytical choices ($\fzero = 7.60$ vs.\ 7.49 Hz, window width, inclusion of noble positions) affect results and are documented in Section 2.4.3
    \item Makes specific, falsifiable predictions testable with any EEG dataset
    \item Connects to fundamental mathematics ($\phisym$) with independent convergence on $\fzero \approx 7.6$ Hz from geophysical monitoring and transient neural events
\end{enumerate}

The golden ratio's unique properties---maximal anti-commensurability enabling within-band independence, and Fibonacci additivity enabling between-band communication---make it optimally suited for neural frequency organization.


%==============================================================================
% References
%==============================================================================

\newpage
\bibliographystyle{apalike}

\begin{thebibliography}{99}

\bibitem[Abelson \& Tukey(1963)]{abelson1963efficient}
Abelson, R. P., \& Tukey, J. W. (1963).
\newblock Efficient utilization of non-numerical information in quantitative analysis: General theory and the case of simple order.
\newblock \textit{The Annals of Mathematical Statistics}, 34(4), 1347--1369.

\bibitem[Arnold(1961)]{arnold1961}
Arnold, V. I. (1961).
\newblock Small denominators. I. Mapping the circle onto itself.
\newblock \textit{Izvestiya Akademii Nauk SSSR. Seriya Matematicheskaya}, 25(1), 21--86.
\newblock (English translation: \textit{American Mathematical Society Translations}, Series 2, 46, 213--284, 1965.)

\bibitem[Badcock et al.(2013)]{badcock2013}
Badcock, N. A., Mousikou, P., Mahajan, Y., de Lissa, P., Thie, J., \& McArthur, G. (2013).
\newblock Validation of the Emotiv EPOC® EEG gaming system for measuring research quality auditory ERPs.
\newblock \textit{PeerJ}, 1, e38.

\bibitem[Berger(1929)]{berger1929}
Berger, H. (1929).
\newblock {\"U}ber das Elektrenkephalogramm des Menschen.
\newblock \textit{Archiv f{\"u}r Psychiatrie und Nervenkrankheiten}, 87, 527--570.

\bibitem[Buzs{\'a}ki(2006)]{buzsaki2006}
Buzs{\'a}ki, G. (2006).
\newblock \textit{Rhythms of the Brain}.
\newblock Oxford University Press.

\bibitem[Debener et al.(2012)]{debener2012}
Debener, S., Minow, F., Emkes, R., Gandras, K., \& De Vos, M. (2012).
\newblock How about taking a low-cost, small, and wireless EEG for a walk?
\newblock \textit{Psychophysiology}, 49(11), 1617--1621.

\bibitem[Donoghue et al.(2020)]{donoghue2020}
Donoghue, T., Haller, M., Peterson, E. J., Varma, P., Sebastian, P., Gao, R., ... \& Voytek, B. (2020).
\newblock Parameterizing neural power spectra into periodic and aperiodic components.
\newblock \textit{Nature Neuroscience}, 23(12), 1655--1665.

\bibitem[Douady \& Couder(1996)]{douady1996}
Douady, S., \& Couder, Y. (1996).
\newblock Phyllotaxis as a dynamical self organizing process.
\newblock \textit{Journal of Theoretical Biology}, 178(3), 255--274.

\bibitem[Goldberger et al.(2002)]{goldberger2002}
Goldberger, A. L., Amaral, L. A., Hausdorff, J. M., Ivanov, P. C., Peng, C. K., \& Stanley, H. E. (2002).
\newblock Fractal dynamics in physiology: alterations with disease and aging.
\newblock \textit{Proceedings of the National Academy of Sciences}, 99(suppl 1), 2466--2472.

\bibitem[Glass \& Mackey(1988)]{glass1988}
Glass, L., \& Mackey, M. C. (1988).
\newblock \textit{From Clocks to Chaos: The Rhythms of Life}.
\newblock Princeton University Press.

\bibitem[Gray \& Singer(1989)]{gray1989}
Gray, C. M., \& Singer, W. (1989).
\newblock Stimulus-specific neuronal oscillations in orientation columns of cat visual cortex.
\newblock \textit{Proceedings of the National Academy of Sciences}, 86(5), 1698--1702.

\bibitem[Iaccarino et al.(2016)]{iaccarino2016}
Iaccarino, H. F., Singer, A. C., Martorell, A. J., Rudenko, A., Gao, F., Gillingham, T. Z., ... \& Tsai, L. H. (2016).
\newblock Gamma frequency entrainment attenuates amyloid load and modifies microglia.
\newblock \textit{Nature}, 540(7632), 230--235.

\bibitem[Krigolson et al.(2017)]{krigolson2017}
Krigolson, O. E., Williams, C. C., Norton, A., Hassall, C. D., \& Colino, F. L. (2017).
\newblock Choosing MUSE: Validation of a low-cost, portable EEG system for ERP research.
\newblock \textit{Frontiers in Neuroscience}, 11, 109.

\bibitem[Livio(2002)]{livio2002}
Livio, M. (2002).
\newblock \textit{The Golden Ratio: The Story of Phi, the World's Most Astonishing Number}.
\newblock Broadway Books.

\bibitem[MacKay(1992)]{mackay1992}
MacKay, R. S. (1992).
\newblock Renormalisation in Area-Preserving Maps.
\newblock \textit{World Scientific}.

\bibitem[Niedermeyer \& Lopes da Silva(2005)]{niedermeyer2005}
Niedermeyer, E., \& Lopes da Silva, F. (2005).
\newblock \textit{Electroencephalography: Basic Principles, Clinical Applications, and Related Fields} (5th ed.).
\newblock Lippincott Williams \& Wilkins.

\bibitem[Nunez \& Srinivasan(2006)]{nunez2006}
Nunez, P. L., \& Srinivasan, R. (2006).
\newblock \textit{Electric Fields of the Brain: The Neurophysics of EEG}.
\newblock Oxford University Press.

\bibitem[Penttonen \& Buzs{\'a}ki(2003)]{penttonen2003}
Penttonen, M., \& Buzs{\'a}ki, G. (2003).
\newblock Natural logarithmic relationship between brain oscillators.
\newblock \textit{Thalamus \& Related Systems}, 2(2), 145--152.

\bibitem[Persinger(2014)]{persinger2014}
Persinger, M. A. (2014).
\newblock Schumann resonance frequencies found within quantitative electroencephalographic activity.
\newblock \textit{International Letters of Chemistry, Physics and Astronomy}, 30, 24--32.

\bibitem[Saroka et al.(2016)]{saroka2016}
Saroka, K. S., Vares, D. E., \& Persinger, M. A. (2016).
\newblock Similar spectral power densities within the Schumann resonance and a large population of quantitative electroencephalographic profiles.
\newblock \textit{PloS One}, 11(1), e0146595.

\bibitem[Schumann(1952)]{schumann1952}
Schumann, W. O. (1952).
\newblock {\"U}ber die strahlungslosen Eigenschwingungen einer leitenden Kugel, die von einer Luftschicht und einer Ionosph{\"a}renh{\"u}lle umgeben ist.
\newblock \textit{Zeitschrift f{\"u}r Naturforschung A}, 7(2), 149--154.

\bibitem[West et al.(1997)]{west1997}
West, G. B., Brown, J. H., \& Enquist, B. J. (1997).
\newblock A general model for the origin of allometric scaling laws in biology.
\newblock \textit{Science}, 276(5309), 122--126.

\bibitem[Lacy(in prep.)]{lacy_sie}
Lacy, M. (in preparation).
\newblock Transient High-Coherence Neural States with Golden Ratio Harmonic Organization at Earth-Resonant Frequencies

\end{thebibliography}


\end{document}
